% concatenated from log-schrodinger.tex
% !TeX spellcheck = en_US
\documentclass[10pt, a4paper, reqno]{amsart}
\usepackage[foot]{amsaddr}
% \makeatletter
% \renewcommand{\email}[2][]{%
% 	\ifx\emails\@empty\relax\else{\g@addto@macro\emails{,\space}}\fi%
% 	\@ifnotempty{#1}{\g@addto@macro\emails{\textrm{#1:}\space}}%
% 	\g@addto@macro\emails{#2}%
% }
% \makeatother
\usepackage[margin=2.5cm]{geometry}
\usepackage{graphicx} 
\usepackage{amsthm, amsmath, amsfonts, amssymb, amscd, enumerate, enumitem, esint, xcolor, url, mathrsfs}
\usepackage{hyperref}

\usepackage[utf8]{inputenc}
\usepackage[T1]{fontenc}
\usepackage{mathtools}
\usepackage{enumerate}
\usepackage{textcomp}

% ---------------------------------------------------
\newtheorem{thm}{Theorem}[section]
\newtheorem{cor}[thm]{Corollary}
\newtheorem{lem}[thm]{Lemma}
\newtheorem{prop}[thm]{Proposition}
\theoremstyle{definition}
\newtheorem{defn}[thm]{Definition}
\theoremstyle{remark}
\newtheorem{rem}[thm]{Remark}
\newtheorem{af}[thm]{Claim}

%---------------------------------------------------------
\DeclareMathOperator{\LLog}{\mathbb{L}\textup{og}}
\DeclareMathOperator{\Log}{\textup{Log}}
\DeclareMathOperator{\supp}{supp}
\DeclareMathOperator*{\esssup}{ess\ sup}
\DeclareMathOperator*{\essinf}{ess\ inf}

\newcommand{\loc}{\textup{loc}}
\newcommand{\glob}{\textup{glob}}
\DeclareMathOperator{\Real}{Re}
\newcommand{\RR}{\mathbb{R}}
\newcommand{\NN}{\mathbb{N}}
\newcommand{\CC}{\mathbb{C}}
\newcommand{\RH}{\textup{RH}}
\newcommand{\Lip}{\textup{Lip}}

\newcommand{\red}[1]{{\color{red!80} #1}}
\newcommand{\blue}[1]{{\color{blue!80} #1}}


%%% ------------------------------------------------------
\numberwithin{equation}{section}
\usepackage[defaultlines=2, all]{nowidow}
\setnoclub[2]
\allowdisplaybreaks

\setlist[enumerate,1]{label=\normalfont{(\alph*)}}

%----------------------------------------------------------


\begin{document}
	
	
	\title[Logarithmic Schr\"odinger operators]
	{Logarithmic Schr\"odinger operators}
	
	
	\author[J. J. Betancor]{Jorge J. Betancor$^1$}
        \address{$^1$Departamento de An\'alisis Matem\'atico, Universidad de La Laguna,\newline
	Campus de Anchieta, Avda. Astrof\'isico S\'anchez, s/n,\newline
	38721 La Laguna (Sta. Cruz de Tenerife), Spain}
	\email[J. J. Betancor]{jbetanco@ull.es}
	
	\author[E. Dalmasso]{Estefan\'ia Dalmasso$^{2}$}
	\address{$^2$Instituto de Matem\'atica Aplicada del Litoral, UNL, CONICET,     FIQ.\newline 
    Colectora Ruta Nac. Nº 168, Paraje El Pozo,\newline S3007ABA, Santa Fe, Argentina}
	\email[E. Dalmasso]{edalmasso@santafe-conicet.gov.ar}

    \author[J. C. Fariña]{Juan Carlos Fariña$^1$}  
        \email[J. C. Fariña]{jcfarina@ull.edu.es}
	
	\author[P. Quijano]{Pablo Quijano$^{2}$}
	\email[P. Quijano]{pquijano@santafe-conicet.gov.ar}
	
	\date{\today}
	\subjclass{35R11, 35J10, 42B37, 47D06}
	
	\keywords{Schr\"odinger, logarthmic operator, nonolocal operator.}

\begin{abstract}
    In this paper we consider the Schr\"odinger operator $\mathcal L_V= -\Delta + V$ in $\mathbb R^d$ with a non negative potential $V$, and $V\not\equiv 0$. We define the logarithmic Schr\"odinger operator $\log \mathcal L_V$ proving its main properties. We obtain a pointwise representation of $\log \mathcal L_V$ when $V$ satisfies a reverse H\"older inequality of exponent $q> \frac{d}{2}$ by using the semigroup of operators $\{T_t^V\}_{t>0}$ generated by $\mathcal L_V$. We consider the Lipschitz function space adapted to the Schr\"odinger setting to solve the initial value problem
    \[\begin{dcases*}
    \frac{\partial u}{\partial t}=-(\log\mathcal L_V)u, & $\text{in } \mathbb R^n\times (0,\infty)$ \\
    u(x,0)=f(x), & $x\in \mathbb R^d$
    \end{dcases*}\]
      in terms of the fractional integral associated with $\mathcal L_V$.
\end{abstract}

\maketitle

\section{Introduction}
Assume that $V$ is a nonnegative potential in $L^1_{\textup{loc}}(\mathbb{R}^d)$. We consider the sesquilinear form $T$ defined by 
\begin{equation*}
    T(f,g) = \int_{\mathbb{R}^d} \nabla f \cdot \overline{\nabla g} \, dx + \int_{\mathbb{R}^d} V f \,\overline{g}\, dx,
\end{equation*}
where $f, g \in D_T = \{h \in L^2(\mathbb{R}^d): \nabla h \in L^2(\mathbb{R}^d) \text{ and } \sqrt{V} h \in L^2(\mathbb{R}^d)\}$. Since $D_T$ is dense in $L^2(\mathbb{R}^d)$ and the quadratic form defined by $T$ is closed and nonnegative, there exists a self-adjoint operator \linebreak $\mathcal{L}_V: D_T =  D(\mathcal{L}_V)  \rightarrow L^2(\mathbb{R}^d)$ such that $\langle \mathcal{L}_V f, g\rangle = T(f,g)$ for all $f, g \in D(\mathcal{L}_V)$ \cite[Theorem VIII.15]{ReSi}. Moreover, for every $f \in C^\infty_c(\mathbb{R}^d)$, the space of $C^\infty$ functions with compact support in $\mathbb{R}^d$, we have $\mathcal{L}_V f = -\Delta f + V f$.

Our objective is to define and study the logarithmic operator $\log \mathcal{L}_V$. This operator is defined using functional calculus. We prove that, alternatively, $\log \mathcal{L}_V$ can be defined as the first variation of the fractional operator $\mathcal{L}_V^{s}$ at $s=0$. We give a pointwise representation of the operator $\log \mathcal{L}_V$. We also solve the initial value problem $\frac{\partial u}{\partial t}=-(\log\mathcal L_V)u$ in $\mathbb R^n\times (0,\infty)$, and $u(x,0)=f(x)$, $x\in \mathbb R^d$, by using fractional integrals and Lipschitz spaces in the Schr\"odinger setting.
%prove that $\log \mathcal{L}_V$, being a nonlocal operator, admits a representation by a local extension problem.

In the last decade, nonlocal operators have been studied by many authors. These operators appear in physics, probability theory, geometry, and applied mathematics (see, for instance,~\cite{CS,  FJW, JX, LL, LPW, Lu, RS, SV}). One of the most important examples of nonlocal operators is the fractional Laplacian $(-\Delta)^s$, $s>0$. The logarithmic Laplacian $\log (-\Delta)$ is a nonlocal operator that has received a lot of attention in recent years (see~\cite{Ch-orderm, ChHW, ChV2, ChV1, ChW, Ch-manifolds, ChX, Feu1, Feu2, HLS, HLW, LW, Lee}). The logarithmic Laplacian $\log (-\Delta)$ is closely connected with the fractional Laplacian $(-\Delta)^s$ by $\log (-\Delta) = \partial_s (-\Delta)^s |_{s=0}$. Recently, the study of $\log(-\Delta)$ has been extended to general Riemannian manifolds (\cite{Ch-manifolds}) and general graphs (\cite{ChX}) %\red{estas citas están en la lista anterior. Son preprints. Revsar antes de enviar si se publicaron}. \blue{El de orden m está publicado con un titulo parecido, ya lo cambié; hay que ver cuando citemos resultados que sean los publicados}


We now define the logarithmic Schr\"odinger operator $\log \mathcal{L}_V$, which is the object of our study. The spectrum $\sigma(\mathcal{L}_V)$ is contained in $[0,+\infty)$. Moreover, $0$ is not an eigenvalue of $\mathcal{L}_V$. Indeed, suppose that $f \in D(\mathcal{L}_V)$ is such that $\mathcal{L}_V f = 0$. Then,
\begin{equation*}
    \langle \mathcal{L}_V f, f \rangle = \int_{\mathbb{R}^d} |\nabla f|^2 \, dx + \int_{\mathbb{R}^d} V |f|^2 \, dx = 0.
\end{equation*}
Hence, $\nabla f = 0$ a.e. in $\mathbb{R}^d$, where the gradient is understood distributionally, and $f(x) = 0$ a.e. \linebreak $x \in \{y \in \mathbb{R}^d : V(y) > 0\}$. Since $f \in L^2(\mathbb{R}^d)$, then $f(x) = 0$ a.e. $x \in \mathbb{R}^d$.

There exists a spectral measure $E_V$ supported on $\sigma(\mathcal{L}_V)$ such that 
\begin{equation*}
    \mathcal{L}_V f = \int_{[0,+\infty)} \lambda \, dE_V(\lambda) f = \int_0^{\infty} \lambda \, dE_V(\lambda) f, \quad f \in D(\mathcal{L}_V), 
\end{equation*}
and 
\begin{equation*}
    D(\mathcal{L}_V) = \left\{ 
    f \in L^2(\mathbb{R}^d) : \int_{0}^\infty \lambda^2 \, d\mu_{f,f}^V(\lambda) < \infty \right\}.
\end{equation*}
Here, for every $f, g \in L^2(\mathbb{R}^d)$, $\mu_{f,g}^V$ is a complex Borel measure given by $\mu_{f,g}^V(U) = \langle E_V(U) f, g \rangle$, for every Borel measurable set $U \subseteq \mathbb{R}$. Note that $E_V(\{0\}) = 0$ because $0$ is not an eigenvalue of $\mathcal{L}_V$.

If $\phi$ is a Borel measurable function on $\mathbb{R}$, we define
\begin{equation*}
    \phi(\mathcal{L}_V) f = \int_0^\infty \phi(\lambda) \, dE_V(\lambda) f, \quad \text{for every } f \in D(\phi(\mathcal{L}_V)),
\end{equation*}
where 
\begin{equation*}
    D(\phi(\mathcal{L}_V)) = \left\{ 
    f \in L^2(\mathbb{R}^d) : \int_{0}^\infty |\phi(\lambda)|^2 \, d\mu_{f,f}^V(\lambda) < \infty \right\}.
\end{equation*}


The logarithmic operator $\log \mathcal{L}_V$ is therefore defined by
\[
(\log \mathcal{L}_V) f = \int_0^\infty \log(\lambda) \, dE_V(\lambda) f, \quad f \in D(\log \mathcal{L}_V),
\]
where
\begin{equation*}
    D(\log \mathcal{L}_V) = \left\{ 
    f \in L^2(\mathbb{R}^d) : \int_{0}^\infty |\log(\lambda)|^2 \, d\mu_{f,f}^V(\lambda) < \infty \right\}.
\end{equation*}

When $V(x) = |x|^2$, $x\in\mathbb{R}^d$, the operator $\mathcal{L}_V$ reduces to the harmonic oscillator operator $H$. In this case, the spectrum $\sigma(H) = \{2k+d\}_{k\in \mathbb{N}_0}$ is discrete. Here $\mathbb{N}_0 = \mathbb{N} \cup \{0\}$. Moreover, for every $k\in \mathbb{N}_0$ and $\alpha = (\alpha_1,\dots,\alpha_d)\in \mathbb{N}_0^d$ such that $\alpha_1+\dots+\alpha_d = k$, 
\begin{equation*}
    H h_\alpha = (2k+d) h_\alpha,
\end{equation*}
where $\displaystyle h_\alpha(x) = \prod_{j=1}^d h_{\alpha_j}(x_j)$, $x=(x_1,\dots,x_d)\in\mathbb{R}^d$, being
\begin{equation*}
    h_\ell (u) = (-1)^\ell (\sqrt{\pi} 2^\ell \ell !)^{-1/2} \frac{ \partial^\ell}{\partial u^\ell} \left(e^{-u^2}\right) e^{u^2/2}, \quad u\in\mathbb{R} \text{ and } \ell\in\mathbb{N}_0.
\end{equation*}
If $\alpha = (\alpha_1,\dots,\alpha_d)\in \mathbb{N}_0^d$, we denote $s(\alpha) = \alpha_1 +\dots + \alpha_d$. We have that
\begin{equation*}
    (\log H) f = \sum_{k\in\mathbb{N}_0} \log(2k+d) \sum_{\substack{\alpha\in \mathbb{N}_0^d \\ s(\alpha)=k}} c_\alpha(f) h_\alpha, \quad \text{for every } f \in D(\log H).
\end{equation*}
Here, 
\begin{equation*}
    D(\log H) =  \left\{ f \in L^2(\mathbb{R}^d): 
    \sum_{k\in\mathbb{N}_0} \sum_{\substack{\alpha\in \mathbb{N}_0^d \\ s(\alpha)=k}} 
    |\log(2k+d)|^2 |c_\alpha(f)|^2 < \infty
    \right\}
\end{equation*}
and, for $\alpha \in \mathbb{N}_0^d$,
\begin{equation*}
    c_\alpha(f) = \int_{\mathbb{R}^d} h_\alpha(x) f(x) \, dx, \quad f \in L^2(\mathbb{R}^d).
\end{equation*}

For $s \in (0,1)$ we define the $s$-power of $\mathcal{L}_V$ by
\begin{equation*}
    \mathcal{L}_V^s f = \int_0^\infty \lambda^s \, dE_V(\lambda) f, \quad f \in D(\mathcal{L}_V^s),
\end{equation*}
where
\begin{equation*}
    D(\mathcal{L}_V^s) = \left\{
    f \in L^2(\mathbb{R}^d) : \int_0^\infty \lambda^{2s} \, d\mu_{f,f}^V(\lambda) < \infty
    \right\}.    
\end{equation*}
Note that $D(\mathcal{L}_V^{s_2}) \subseteq D(\mathcal{L}_V^{s_1})$ 
provided that $0 < s_1 \leq s_2 < 1$. If $f \in D(\mathcal{L}_V^\beta)$ for some $\beta \in (0,1)$, we have that
\begin{equation*}
    \lim_{s \rightarrow 0^+} \mathcal{L}_V^s f = f \quad \text{in } L^2(\mathbb{R}^d). 
\end{equation*}

   We define the operator $\LLog \mathcal{L}_V$ as follows:
\begin{equation*}
    (\LLog \mathcal{L}_V) f = \frac{d}{ds}(\mathcal{L}_V^s f)\big|_{s=0} = \lim_{s\rightarrow 0^+} \frac{\mathcal{L}_V^s f - f}{s},
\end{equation*}
provided that 
\begin{equation*}
    f \in D(\LLog \mathcal{L}_V) = \left\{ f \in L^2(\mathbb{R}^d) : f \in D(\mathcal{L}_V^\beta) \text{ for some } \beta \in (0,1) \text{ and } \lim_{s\rightarrow 0^+} \frac{\mathcal{L}_V^s f - f}{s} \text{ exists in } L^2(\mathbb{R}^d) \right\}.
\end{equation*}
We define, for every $t > 0$, the operator $T_t^V$ by 
\begin{equation*}
    T_t^V f = \int_0^\infty e^{-\lambda t} \, dE_V(\lambda) f, \quad f \in L^2(\mathbb{R}^d).
\end{equation*}

The family $\{T_t^V\}_{t>0}$ is a $C_0$-semigroup of operators on $L^2(\mathbb{R}^d)$. The operator $\mathcal{L}_V$ is the infinitesimal generator of $\{T_t^V\}_{t>0}$. We can write, for every $t > 0$, $T_t^V = \phi_t(\mathcal{L}_V)$, where $\phi_t(\lambda) = e^{-\lambda t}$, $\lambda \in (0,\infty)$.

The Frullani integral representation (see~\cite[p. 363]{GR}) allows us to write
\begin{equation*}
    \log(\lambda) = \int_0^\infty \frac{e^{-t} - e^{-\lambda t}}{t} \, dt, \quad \lambda > 0. 
\end{equation*}
This formula motivates the following definition:
\begin{equation*}
    (\Log \mathcal{L}_V) f = \lim_{m \rightarrow \infty} \int_{1/m}^m \frac{e^{-t}I - T_t^V}{t} f \, dt, \quad f \in D(\Log \mathcal{L}_V),
\end{equation*}
where the integrals are understood in the $L^2(\mathbb{R}^d)$-Bochner sense and the limit is in $L^2(\mathbb{R}^d)$. Here
$D(\Log \mathcal{L}_V)$ is the set of all functions $f \in L^2(\mathbb{R}^d)$ such that $\displaystyle \frac{e^{-t} I - T_t^V}{t} f$ is $L^2(\mathbb{R}^d)$-Bochner integrable on $(1/m,m)$ for every $m \in \mathbb{N}$ and the limit
\begin{equation*}
    \lim_{m \rightarrow \infty} \int_{1/m}^m \frac{e^{-t} I - T_t^V}{t} f \, dt
\end{equation*}
exists in $L^2(\mathbb{R}^d)$. We recall that $\displaystyle \frac{e^{-t} I - T_t^V}{t} f$ is $L^2(\mathbb{R}^d)$-Bochner integrable on $(1/m,m)$ for $m \in \mathbb{N}$ when it is $L^2(\mathbb{R}^d)$-strongly measurable on $((1/m,m), dt)$ and 
\begin{equation*}
    \int_{1/m}^m \left\| \frac{e^{-t} I - T_t^V}{t} f \right\|_{L^2(\mathbb{R}^d)} \, dt < \infty.
\end{equation*}

The following theorem establishes the equivalence between the different definitions of the logarithmic operator.

\begin{thm}\label{thm: 1.1}
    Let $f \in L^2(\mathbb{R}^d)$. We have that
    \begin{enumerate}
        \item\label{item: teo 1.1 - a} $f \in D(\mathcal{L}_V^{s_0}) \cap D(\log \mathcal{L}_V)$ for some $s_0 \in (0,1)$, if and only if $f \in D(\LLog \mathcal{L}_V)$ and $(\log \mathcal{L}_V) f = (\LLog \mathcal{L}_V) f$.
        \item\label{item: teo 1.1 - b} If $f \in D(\log \mathcal{L}_V)$, then $f \in D(\Log \mathcal{L}_V)$ and $(\log \mathcal{L}_V) f = (\Log \mathcal{L}_V) f$.
    \end{enumerate}
\end{thm}

Our next objective is to obtain a pointwise representation of $\log \mathcal{L}_V$ on $C_c^\infty(\mathbb{R}^d)$. Note that, as will be established in Proposition~\ref{prop: 2.1}, we have $C_c^\infty(\mathbb{R}^d) \subseteq D(\log \mathcal{L}_V)$. 

A key difference from the classical case arises from the semigroup properties: unlike the classical heat semigroup $\{W_t\}_{t>0}$ generated by $-\Delta$, which satisfies $W_t 1 = 1$ for all $t > 0$ (the Markovian property), the semigroup $\{T_t^V\}_{t>0}$ is not Markovian, that is, $T_t^V 1 \neq 1$ for $t > 0$. This fundamental difference produces significant distinctions between the pointwise representation for $\log(-\Delta)$ established in~\cite[Theorem 1.1]{ChW} and the one we shall derive for $\log(\mathcal{L}_V)$.

In order to establish our pointwise representation of $\log \mathcal{L}_V$, we consider a special class of potentials satisfying reverse H\"older inequalities. For $1 < q < \infty$, we say that a locally integrable function $w$ is in the $q$-reverse H\"older class $\RH_q$ when there exists $C > 0$ such that
\begin{equation*}
    \left(\frac{1}{|B|} \int_B w^q \, dx\right)^{1/q} \leq \frac{C}{|B|} \int_B w \, dx,
\end{equation*}
for every ball $B$ in $\mathbb{R}^d$. Functions in $\RH_q$ classes were first studied by Gehring (\cite{G}) and Coifman--Fefferman (\cite{CF}). The condition $V \in \RH_q$ for some $1 < q < \infty$ naturally appears in the study of harmonic analysis operators in the Schr\"odinger setting, such as Riesz transforms, Littlewood--Paley functions, spectral multipliers, commutators, and others (see~\cite{BFHR, BHQ, Dz1, Dz2, DGMTZ, GHY, Shen95, ZT},
and the references therein). 

Suppose that $d \geq 3$ and $V \in \RH_q$ with $q > d/2$. As in~\cite[Eq.~(0.14)]{Shen95}, we define the function $\rho: \mathbb{R}^d \rightarrow (0,\infty)$ by
\begin{equation*}
    \rho(x) = \sup\left\{ r > 0 : \frac{1}{r^{d-2}} \int_{B(x,r)} V(y) \, dy \leq 1 \right\}.
\end{equation*}

The function $\rho$ was introduced in~\cite{Shen94} for potentials $V$ satisfying, for some $C > 0$, 
\begin{equation*}
    \max_{x \in B} V(x) \leq \frac{C}{|B|} \int_B V(x) \, dx,
\end{equation*}
for every ball $B$ in $\mathbb{R}^d$, in the context of studying the Neumann problem for $\mathcal{L}_V$. This auxiliary function $\rho$ plays a fundamental role in harmonic analysis in the Schr\"odinger setting.

It is also proved in~\cite{Shen94} that this function $\rho$ satisfies the following estimates: \begin{equation}\label{eq: rox_vs_roy} c_\rho^{-1} \, \rho(x)\left(1+\frac{|x-y|}{\rho(x)}\right)^{-N_0} \ \leq \ \rho(y) \ \leq \ c_\rho\,\rho(x)\left(1+\frac{|x-y|}{\rho(x)}\right)^{\frac{N_0}{N_0+1}}, \end{equation} for some positive constants $c_\rho$ and $N_0$ independent of $x$ and $y$.
As a consequence of \eqref{eq: rox_vs_roy}, we have in particular that $\rho(x)\sim \rho(y)$ whenever $|x-y|\leq \rho(x)$. That is, there exists a constant $C>0$, independent of $x$ and $y$, such that
$
C^{-1}\,\rho(x) \leq \rho(y) \leq C\,\rho(x)
$ whenever $|x-y|\leq \rho(x)$.

For every $t > 0$ there exists a measurable function $T_t^V : \mathbb{R}^d \times \mathbb{R}^d \rightarrow (0,\infty)$ such that
\begin{equation}\label{eq: 1.1}
    T_t^V(f)(x) = \int_{\mathbb{R}^d} T_t^V(x,y) f(y) \, dy,
\end{equation}
for every $f \in L^2(\mathbb{R}^d)$. Moreover,~\eqref{eq: 1.1} allows us to define a semigroup of contractions in $L^p(\mathbb{R}^d)$, ${1 \leq  p < \infty}$. Since $V$ is nonnegative, the Feynman--Kac formula implies that
\begin{equation}\label{eq: 1.2}
    0 \leq T_t^V(x,y) \leq T_t(x-y), \quad x, y \in \mathbb{R}^d \text{ and } t > 0,
\end{equation}
where $T_t(z) = (4\pi t)^{-d/2} e^{-|z|^2/(4t)}$, $z \in \mathbb{R}^d$ and $t > 0$, is the classical heat kernel.

Estimate~\eqref{eq: 1.2} can be improved using that $V \in \RH_q$ with $q > d/2$. As stated in \cite[Proposition~2]{DGMTZ} (see also \cite{DzZ2, Kur}), for every $N \in \mathbb{N}$ there exists $C_N> 0$ such that 
\begin{equation}\label{eq: 1.3}
    T_t^V(x,y) \leq \frac{C_N}{t^{d/2}} e^{-\frac{|x-y|^2}{5t}} \left( 1 + \frac{\sqrt{t}}{\rho(x)} + \frac{\sqrt{t}}{\rho(y)}\right)^{-N}, 
\end{equation}
for all $x, y \in \mathbb{R}^d$ and $t > 0$. Also,  \cite[Proposition~4.11]{DzZ2} provides a regularity estimate: there exist $\eta, c > 0$ such that for every $N \in \mathbb{N}$ there exists a constant $C_N > 0$ for which
\begin{equation}\label{eq: 1.4}
    |T_t^V(x+h,y) - T_t^V(x,y)| \leq \frac{C_N}{t^{d/2}} \left(\frac{|h|}{\sqrt{t}} \right)^{\eta} e^{-c\frac{|x-y|^2}{t}} \left( 1 + \frac{\sqrt{t}}{\rho(x)} + \frac{\sqrt{t}}{\rho(y)}\right)^{-N}, 
\end{equation}
whenever $|h| < \sqrt{t}$, $x, y \in \mathbb{R}^d$ and $t > 0$.

The operator $\mathcal{L}_V$ can be seen as a nice perturbation of the Euclidean Laplacian $-\Delta$ in the following sense: there exists a nonnegative Schwartz class function $\omega$ on $\mathbb{R}^d$ and $\delta > 0$ such that 
\begin{equation}\label{eq: 1.5}
    |T_t^V(x,y) - T_t(x-y)| \leq \begin{cases}
\left(\dfrac{\sqrt{t}}{\rho(x)}\right)^\delta \omega_t(x-y), & \text{if } \sqrt{t} \leq \rho(x),  \\
\omega_t(x-y), & \text{if } \sqrt{t} > \rho(x),
\end{cases}
\end{equation}
for $x, y \in \mathbb{R}^d$, where $\omega_t(z) = t^{-d/2} \omega(z/\sqrt{t})$, $z \in \mathbb{R}^d$ and $t > 0$ (see \cite[Proposition~2.16]{DzZ3}).
We now establish our pointwise representation for $\log(\mathcal{L}_V)$. This result extends, in particular, \cite[Theorem~1.1]{Feu1} where $V \equiv 1$ is considered. 

\begin{thm}\label{thm: 1.2} 
    Let $d \geq 3$ and $q > d/2$. Suppose that $V \in \RH_q$ and $f \in C^\infty_c(\mathbb{R}^d)$. 
    \begin{enumerate}
        \item \label{item: teo 1.2 - a} We have that
        \begin{equation}\label{eq: teo 1.2}
        \begin{split}
            (\log\mathcal{L}_V)f(x) & = -\int_{B(x,1)} (f(y)-f(x)) \int_0^\infty \frac{T_t^V(x,y)}{t} \, dt \, dy \\
            & \quad - \int_{\mathbb{R}^d \setminus B(x,1)} f(y) \int_0^\infty \frac{T_t^V(x,y)}{t} \, dt \, dy - K(x)f(x),
        \end{split}
        \end{equation}
        for almost every $x \in \mathbb{R}^d$, where
        \begin{equation*}
        \begin{split}
            K(x) & = 2 \log \rho(x) + \int_0^{\rho^2(x)} \int_{\mathbb{R}^d} \frac{T_t^V(x,y) - T_t(x-y)}{t} \, dy \, dt \\
            & \quad - \int_0^{\rho^2(x)} \int_{\mathbb{R}^d \setminus B(x,1)} \frac{T_t^V(x,y)}{t} \, dy \, dt \\
            & \quad + \int_{\rho^2(x)}^\infty \int_{B(x,1)} \frac{T_t^V(x,y)}{t} \, dy \, dt +\gamma,
        \end{split}
        \end{equation*}
        and $\gamma$ denotes the Euler-Mascheroni constant.

        \item \label{item: teo 1.2 - b} For every $1 < p \leq \infty$,
        \begin{equation*}
            \lim_{s \rightarrow 0^{+}} \frac{\mathcal{L}_V^s f - f}{s} = (\log \mathcal{L}_V) f, 
        \end{equation*}
        in $L^p(\mathbb{R}^d)$.
    \end{enumerate}
\end{thm}

Note that the function $K$ is not constant. This reflects precisely the difference mentioned earlier: the semigroup $\{T_t^V\}_{t>0}$ is not Markovian, unlike the classical heat semigroup. Consequently, the constant term in the pointwise representation of $\log \mathcal{L}_V$ becomes a non-trivial function $K(x)$ depending on the auxiliary function $\rho(x)$, which encodes information about the potential $V$. This fact, combined with the lack of an explicit formula for the heat kernel $T_t^V(x,y)$---we only have the estimates \eqref{eq: 1.2}, \eqref{eq: 1.3}, \eqref{eq: 1.4} and \eqref{eq: 1.5}---makes the proof of Theorem~\ref{thm: 1.2} substantially more intricate than that of \cite[Theorem~1.1]{ChW} for the Euclidean case.
A small modification of the proof of Theorem~\ref{thm: 1.2} allows us to see that if $f \in \Lip^\theta(\mathbb{R}^d)$, the $\theta$-Lipschitz space on $\mathbb{R}^d$, for some $\theta \in (0,1]$, and 
\begin{equation*}
    \int_{\mathbb{R}^d} \frac{|f(y)|}{(1+|y|)^{d}} \, dy < \infty,
\end{equation*}
then
\begin{equation*}
\begin{split}
    \lim_{s \rightarrow 0^+} \frac{1}{s} \left( \frac{1}{\Gamma(-s)} \int_0^\infty \frac{T_t^V f(x) - f(x)}{t^{s+1}} \, dt - f(x) \right) & = - \int_{B(x,1)} (f(y)-f(x)) \int_0^\infty \frac{T_t^V(x,y)}{t} \, dt \, dy \\
    & \quad -\int_{\mathbb{R}^d \setminus B(x,1)} f(y) \int_0^\infty \frac{T_t^V(x,y)}{t} \, dt \, dy - K(x)f(x),
\end{split}
\end{equation*}
for $x \in \mathbb{R}^d$. Thus, the domain of $\log \mathcal{L}_V$ can be extended according to the last identity.

Suppose now $V(x) = 1$, $x \in \mathbb{R}^d$. In this case $T_t^V = e^{-t} T_t$, $t > 0$. By manipulating formula~\eqref{eq: teo 1.2} we get for $f$ as in Theorem~\ref{thm: 1.2}:
\begin{equation*}
    \begin{split}
        (\log \mathcal{L}_V) f(x) & = -\int_{\mathbb{R}^d} (f(y)-f(x)) \int_0^\infty \frac{e^{-t}T_t(x-y)}{t} \, dt \, dy \\
        & \quad - f(x) \left( 2\log\rho(x) + \int_0^{\rho^2(x)} \frac{e^{-t} - 1}{t} \, dt + \int_{\rho^2(x)}^\infty \frac{e^{-t}}{t} \, dt - \Gamma'(1) \right).
    \end{split}
\end{equation*}
Since
\begin{equation*}
    \log(z) + \int_0^z \frac{e^{-t} - 1}{t} \, dt + \int_z^\infty \frac{e^{-t}}{t} \, dt = -\gamma \quad \text{for } z \in (0,\infty),
\end{equation*}
and $\Gamma'(1) = -\gamma$, we conclude that 
\begin{equation*}
    (\log \mathcal{L}_V) f(x) = -\int_{\mathbb{R}^d} (f(y)-f(x)) \int_0^\infty \frac{e^{-t}}{t} T_t(x-y) \, dt \, dy, \quad \text{a.e. } x \in \mathbb{R}^d.
\end{equation*}

On the other hand, we have that
\begin{equation*}
    \int_0^\infty \frac{e^{-t}}{t} T_t(x-y) \, dt = \int_0^\infty \frac{e^{-t - \frac{|x-y|^2}{4t}}}{(4\pi t)^{d/2} t} \, dt = \frac{2}{(4\pi)^{d/2}} \left(\frac{|x-y|}{2}\right)^{-d/2} \mathcal{K}_{d/2}(|x-y|), \quad x, y \in \mathbb{R}^d,
\end{equation*}
where $\mathcal{K}_\nu$ denotes the modified Bessel function of the second kind of order $\nu$. Thus, \cite[Theorem~1.1]{Feu1} is a special case of Theorem~\ref{thm: 1.2} here.


Fall and Felli (\cite{FaFe2} and \cite{FaFe1}) studied the essential self-adjointness and established unique continuation properties for a relativistic Schr\"odinger operator with a singular homogeneous potential defined by 
\begin{equation*}
    H = (-\Delta + m^2)^s - \frac{a(x/|x|)}{|x|^{2s}} - h(x), 
\end{equation*}
where $m \geq 0$ and $s \in (0,1)$. In~\cite[Eq.~(1.3)]{FaFe1} the authors give a pointwise representation for the operator $(-\Delta + m^2)^s$. A pointwise representation for the operator $\log(-\Delta + m^2)$ appears here as a special case of Theorem~\ref{thm: 1.2}.

In~\cite{BHS} and~\cite{MSTZ}, Lipschitz spaces in the Schr\"odinger setting were introduced as follows: we say that a function $f \in \Lip_V^\alpha$ with $0 < \alpha < 1$ when 
\begin{equation*}
    \|\rho^{-\alpha} f\|_\infty < \infty \quad \text{and} \quad \sup_{|h|>0} \frac{\|f(\cdot + h) - f(\cdot)\|_\infty}{|h|} < \infty .
\end{equation*}
The definition of $\Lip_V^\alpha$ was extended to $0 < \alpha < 2$ in~\cite{dLT}.

For every $\alpha > 0$ we consider the negative fractional power $\mathcal{L}_V^{-\alpha}$ of $\mathcal{L}_V$ defined by
\begin{equation*}
    \mathcal{L}_V^{-\alpha} f = \int_0^\infty \lambda^{-\alpha} \, dE_V(\lambda) f, \quad f \in D(\mathcal{L}_V^{-\alpha}), 
\end{equation*}
where
\begin{equation*}
    D(\mathcal{L}_V^{-\alpha}) = \left\{ f \in L^2(\mathbb{R}^d) : \int_0^\infty \lambda^{-2\alpha} \, d\mu_{f,f}^V(\lambda) < \infty \right\}.
\end{equation*}
We are going to see in Proposition~\ref{propo: 2.5} that $\mathcal{L}_V^{-\alpha}$ can be expressed in terms of the semigroup $\{T_t^V\}_{t>0}$ as 
\begin{equation}\label{eq: 1.6}
    \mathcal{L}_V^{-\alpha} f = \frac{1}{\Gamma(\alpha)} \int_0^\infty T_t^V(f) t^{\alpha-1} \, dt,
\end{equation}
By using~\eqref{eq: 1.1} we can see that the integral in~\eqref{eq: 1.6} allows us to extend $\mathcal{L}_V^{-\alpha}$ to $L^p(\mathbb{R}^d)$ as a bounded operator from $L^p(\mathbb{R}^d)$ into $L^q(\mathbb{R}^d)$ where $1/q = 1/p - 2\alpha/d$, $1 < p < d/(2\alpha)$, and from $L^1(\mathbb{R}^d)$ into $L^{d/(d-2\alpha), \infty}(\mathbb{R}^d)$. As can be seen in equation~\eqref{eq: 1.2}, the integral kernel of the operator $T_t^V$ has a better behavior far away from the diagonal than the classical heat kernel $T_t$. In~\cite{DGMTZ} it was proved that $\mathcal{L}_V^{-\alpha}$ maps $L^{d/(2\alpha)}(\mathbb{R}^d)$ into a $\textup{BMO}$-space denoted by $\textup{BMO}_{\mathcal{L}_V}$ adapted to the Schr\"odinger setting. In \cite[Theorem 1.2]{MSTZ} (see also~\cite[Theorem 1.6]{dLT}) the authors proved that $\mathcal{L}_V^{-\alpha}$ maps $\Lip_V^\theta$ into $\Lip_V^{\theta - 2\alpha}$ when $0 < \theta < 1$, $\alpha > 0$ and $\theta + 2\alpha < 1$.

Suppose now that $0$ does not belong to the spectrum $\sigma(\mathcal{L}_V)$. Then, there exists $a > 0$ such that $\sigma(\mathcal{L}_V) \subseteq [a, +\infty)$. Since the function $\phi(\lambda) = \lambda^{-\alpha}$, $\lambda \in (0,+\infty)$ is bounded on $[a,+\infty)$, $\mathcal{L}_V^{-\alpha}$ is a bounded operator on $L^2(\mathbb{R}^d)$, for every $\alpha > 0$. Moreover, $\mathcal{L}_V^{-\alpha} \circ \mathcal{L}_V^{-\beta} = \mathcal{L}_V^{-(\alpha+\beta)}$, for every $\alpha, \beta > 0$, and $\lim_{\alpha \rightarrow 0^+} \mathcal{L}_V^{-\alpha} f = f$ in $L^2(\mathbb{R}^d)$ for every $f \in L^2(\mathbb{R}^d)$.

The family $\{\mathcal{L}_V^{-\alpha}\}_{\alpha>0}$ is a $C_0$-semigroup of operators on $L^2(\mathbb{R}^d)$. We will prove in Proposition~\ref{propo: 2.3} that the infinitesimal generator of $\{\mathcal{L}_V^{-\alpha}\}_{\alpha>0}$ is the logarithmic operator $\log \mathcal{L}_V$. We recall here that $\sigma(H) = \{2k+d\}_{k\in \mathbb{N}_0}$ when $H$ represents the harmonic oscillator operator.

Our final objective is to study the following Cauchy problem

\begin{equation}\label{eq: 1.7}
    \begin{cases}
        \dfrac{\partial u}{\partial t} (x,t) + (\log \mathcal{L}_V) (u(\cdot,t))(x) = 0, & x \in \mathbb{R}^d, \; t > 0,\\
        u(x,0) = f(x), & x \in \mathbb{R}^d,
    \end{cases}
\end{equation}
in $\mathbb{R}^d \times (0,\infty)$. If $0 \notin \sigma(\mathcal{L}_V)$, then, for every $f \in L^2(\mathbb{R}^d)$, $u(x,t) = \mathcal{L}_V^{-t} (f)(x)$ is an $L^2$-solution of~\eqref{eq: 1.7}.

\begin{thm}\label{thm: 1.3}
    Let $d \geq 3$ and $V \in \RH_q$ with $q > d/2$. Suppose that $f \in L^1(\mathbb{R}^d) \cap \Lip_V^\theta \cap C^1(\mathbb{R}^d)$, for some $\theta \in (0,1)$, and $\nabla f \in L^1(\mathbb{R}^d)$. Then, the function $u$ defined by
    \begin{equation*}
        u(x,t) = \frac{1}{\Gamma(t)} \int_0^\infty T_z^V (f)(x) z^{t-1} \, dz, \quad x \in \mathbb{R}^d \text{ and } 0 < t < 1 - \theta,
    \end{equation*}
    satisfies the following properties:
    \begin{enumerate}
        \item $u$ and $\dfrac{\partial u}{\partial t}$ are continuous functions in $\mathbb{R}^d \times (0,1-\theta)$. 
        \item $\dfrac{\partial u}{\partial t}(x,t) = -(\log \mathcal{L}_V)(u(\cdot,t))(x)$, for $(x,t) \in \mathbb{R}^d \times (0,1-\theta)$.  
        \item For every $x \in \mathbb{R}^d$, $\displaystyle \lim_{t\rightarrow 0^+} u(x,t) = f(x)$.
    \end{enumerate}
\end{thm}

We prove Theorem~\ref{thm: 1.1} in Section~\ref{sec: 2}, Theorem~\ref{thm: 1.2} in Section~\ref{sec: 3} and Theorem~\ref{thm: 1.3} in Section~\ref{sec: 4}. Throughout this paper $C$ and $c$ always represent positive constants that may change in each occurrence.

\section{About the logarithmic operator of \texorpdfstring{$\mathcal{L}_V$}{LV}}\label{sec: 2}


In this section we prove Theorem~\ref{thm: 1.1} together with some properties of the logarithmic operator of $\mathcal{L}_V$. As we mentioned before, we define 
\[
(\log \mathcal{L}_V) f = \int_0^\infty \log(\lambda) \, dE_V(\lambda) f, \quad f \in D(\log \mathcal{L}_V),
\]
where
\begin{equation*}
    D(\log \mathcal{L}_V) = \left\{ 
    f \in L^2(\mathbb{R}^d) : \int_{0}^\infty |\log(\lambda)|^2 \, d\mu_{f,f}^V(\lambda) < \infty \right\}.
\end{equation*}

Here $E_V$ denotes the spectral measure of $\mathcal{L}_V$ and, for every $f \in L^2(\mathbb{R}^d)$, $\mu_{f,f}^V$ is the Borel measure defined by $\mu_{f,f}^V(A) = \langle E_V(A) f, f \rangle$ for each Borel set $A \subseteq (0,\infty)$.

\begin{prop}\label{prop: 2.1}
    Let $d \geq 3$ and $V \in \RH_q$ with $q > d/2$. Then $D(\mathcal{L}_V^{-\beta}) \cap D(\mathcal{L}_V^{\alpha}) \subseteq D(\log \mathcal{L}_V)$, for every $\alpha, \beta \in (0,1)$, and $C_c^\infty (\mathbb{R}^d) \subseteq D(\log \mathcal{L}_V)$.
\end{prop}
\begin{proof}
    Let $f \in L^2(\mathbb{R}^d)$ and $\alpha, \beta \in (0,1)$. We can write
    \begin{equation*}
        \begin{split}
            \int_0^\infty (\log \lambda)^2 \, d\mu_{f,f}^V(\lambda) & \leq C 
            \left( \int_0^1 \lambda^{-2\beta} \, d\mu_{f,f}^V(\lambda) + \int_1^\infty \lambda^{2\alpha} \, d\mu_{f,f}^V(\lambda)\right) \\
            & \leq C 
            \left( \int_0^\infty \lambda^{-2\beta} \, d\mu_{f,f}^V(\lambda) + \int_0^\infty \lambda^{2\alpha} \, d\mu_{f,f}^V(\lambda)\right).
        \end{split}
    \end{equation*}
    We deduce that $f \in D(\log \mathcal{L}_V)$ provided that $f \in D(\mathcal{L}_V^{-\beta}) \cap D(\mathcal{L}_V^{\alpha})$.

    On the other hand, $D(\mathcal{L}_V) \subseteq D(\mathcal{L}_V^\alpha)$. Hence, $C_c^\infty(\mathbb{R}^d) \subseteq D(\mathcal{L}_V^\alpha)$. Moreover, $\mathcal{L}_V^{-\beta}$ is bounded from $L^p(\mathbb{R}^d)$ into $L^q(\mathbb{R}^d)$ when $1/q = 1/p - (2\beta)/d$ and $1 < p < d/(2\beta)$. In particular, $\mathcal{L}_V^{-\beta}$ is bounded from $L^{(2d)/(d+2\beta)}(\mathbb{R}^d)$ into $L^2(\mathbb{R}^d)$, so $L^{(2d)/(d+2\beta)}(\mathbb{R}^d) \subseteq D(\mathcal{L}_V^{-\beta})$. Since $C_c^\infty(\mathbb{R}^d) \subseteq L^{(2d)/(d+2\beta)}(\mathbb{R}^d)$, we conclude that $C_c^\infty(\mathbb{R}^d) \subseteq D(\mathcal{L}_V^{-\beta}) \cap D(\mathcal{L}_V^\alpha) \subseteq D(\log \mathcal{L}_V)$.
\end{proof}

We define, for every $\beta \in \mathbb{R}$, $J_\beta = \mathcal{L}_V^{i\beta}=\phi_\beta(\mathcal{L}_V)$, where $\phi_\beta(\lambda)=\lambda^{i\beta}$, $\lambda\in (0,\infty)$. Since this function is bounded for $\lambda \in (0,\infty)$, the operator $J_\beta$ is bounded on $L^2(\mathbb{R}^d)$ for every $\beta \in \mathbb{R}$. We have that 
\[
J_\alpha \circ J_\beta = J_{\alpha + \beta}, \quad \alpha, \beta \in \mathbb{R}.
\]
Hence $\{J_\beta\}_{\beta \in \mathbb{R}}$ is a group of bounded operators on $L^2(\mathbb{R}^d)$.
\begin{prop}\label{propo: 2.2}
    The operator $i\log \mathcal{L}_V$ is the infinitesimal generator of the group $\{J_\beta\}_{\beta \in \mathbb{R}}$ in $L^2(\mathbb{R}^d)$. 
\end{prop}
\begin{proof}
    We denote by $G$ the infinitesimal generator of $\{J_\beta\}_{\beta \in \mathbb{R}}$ in $L^2(\mathbb{R}^d)$, that is, 
    \[
    Gf = \lim_{h \rightarrow 0} \frac{J_h f - f}{h}, \quad \text{in } L^2(\mathbb{R}^d),
    \]
    for every $\displaystyle f \in D(G) = \left\{ g \in L^2(\mathbb{R}^d) : \lim_{h \rightarrow 0} \frac{J_h g - g}{h} \text{ exists in } L^2(\mathbb{R}^d) \right\}$.

    Let $f \in D(\log \mathcal{L}_V)$. We are going to show that
    \[
    \lim_{h \rightarrow 0} \frac{J_h f - f}{h} = i(\log \mathcal{L}_V) f, \quad \text{in } L^2(\mathbb{R}^d),
    \]
    or equivalently that
    \[
    \left\| \frac{J_h f - f}{h} - i(\log \mathcal{L}_V) f \right\|_{L^2(\mathbb{R}^d)}^2 = 
    \int_0^\infty \left| \frac{\lambda^{ih}-1}{h} - i\log(\lambda)\right|^2 \, d\mu_{f,f}^V(\lambda) \rightarrow 0 
    \]
    as $h \rightarrow 0$. 

    We note that $\left| \frac{\lambda^{ih}-1}{h}\right| \leq |\log \lambda|$ for all $h \in \mathbb{R}$ and $\lambda > 0$. Then, 
    \[
    \left| \frac{\lambda^{ih}-1}{h} - i\log\lambda \right| \leq 2 |\log\lambda|. 
    \]

    Since $f \in D(\log \mathcal{L}_V)$, the Dominated Convergence Theorem leads to 
    \[
    \lim_{h \rightarrow 0} \int_0^\infty \left| \frac{\lambda^{ih}-1}{h} - i\log\lambda \right|^2 \, d\mu_{f,f}^V(\lambda) = 0. 
    \]
    We conclude that $f \in D(G)$ and $Gf = i(\log\mathcal{L}_V) f$.

    Conversely, suppose that $f \in D(G)$. Then there exists $g \in L^2(\mathbb{R}^d)$ such that $Gf = g$, that is,
    \[
    \lim_{h \rightarrow 0} \left\| \int_0^\infty \frac{\lambda^{hi}-1}{h} \, dE_V(\lambda) f - g \right\|_{L^2(\mathbb{R}^d)} = 0.
    \]

    Let $(h_k)_{k \in \mathbb{N}} \subseteq (0,\infty)$ be a sequence converging to zero. The sequence
    \[
    \left\{ \int_0^\infty \frac{\lambda^{h_k i} - 1}{h_k} \, dE_V(\lambda) f \right\}_{k \in \mathbb{N}}
    \]
    is Cauchy in $L^2(\mathbb{R}^d)$, or equivalently, the sequence $\displaystyle \left\{ \frac{\lambda^{h_k i}-1}{h_k}\right\}_{k \in \mathbb{N}}$ is Cauchy in $L^2((0,\infty), d\mu_{f,f}^V)$. Hence, there exists $h \in L^2((0,\infty), d\mu_{f,f}^V)$ such that
    \[
    \lim_{k \rightarrow \infty} \frac{\lambda^{h_k i}-1}{h_k} = h(\lambda) \quad \text{in } L^2((0,\infty), d\mu_{f,f}^V).
    \]
    On the other hand, for each $\lambda > 0$,
    \[
    \lim_{k \rightarrow \infty} \frac{\lambda^{h_k i}-1}{h_k} = i\log\lambda.
    \]
    Therefore, $h(\lambda) = i\log\lambda$ for $\mu_{f,f}^V$-almost every $\lambda > 0$.

    We have proved that 
    \[ 
    \left\| \int_0^\infty \frac{\lambda^{h_k i}-1}{h_k} \, dE_V(\lambda) f - \int_0^\infty i\log\lambda \, dE_V(\lambda) f \right\|_{L^2(\mathbb{R}^d)}^2 = 
    \int_0^\infty \left| \frac{\lambda^{h_k i}-1}{h_k} - i\log\lambda \right|^2 \, d\mu_{f,f}^V(\lambda) \rightarrow 0
    \]
    as $k \rightarrow \infty$. The arbitrariness of the sequence $(h_k)_{k \in \mathbb{N}}$ allows us to conclude that $f \in D(\log \mathcal{L}_V)$ and $g = i(\log \mathcal{L}_V) f$.
\end{proof}

By proceeding in a similar way we can prove the following property:

\begin{prop}\label{propo: 2.3}
    Suppose that $0 \notin \sigma(\mathcal{L}_V)$. Then $-\log \mathcal{L}_V$ is the infinitesimal generator of the $C_0$-semigroup $\{\mathcal{L}_V^{-t}\}_{t>0}$ of operators on $L^2(\mathbb{R}^d)$.
\end{prop}

We proceed now with the proof of Theorem~\ref{thm: 1.1}.

\begin{proof}[Proof of Theorem~\ref{thm: 1.1}]
We begin proving part~\ref{item: teo 1.1 - a}. Suppose that $f \in D(\mathcal{L}_V^{s_0}) \cap D(\log \mathcal{L}_V)$ for some $s_0 \in (0,1)$. Then $f \in D(\mathcal{L}_V^s)$ for every $0 < s < s_0$. We have that
\[
\frac{|\lambda^s-1|}{s} \leq C \begin{cases}
    |\log \lambda|, & \text{if } 0 < \lambda < 1, \\
    |\log \lambda| \lambda^s \leq C \lambda^{s_0}, & \text{if } 0 < s < s_0 \text{ and } \lambda \geq 1.
\end{cases}
\]

Then, 
\[ 
\left| \frac{\lambda^{s}-1}{s} - \log \lambda\right|^2 \leq C (\lambda^{2s_0} + |\log \lambda|^2)
\]
for $\lambda > 0$ and $0 < s < s_0$. Since $\int_0^\infty \lambda^{2s_0} \, d\mu_{f,f}^V(\lambda) < \infty$ and $\int_0^\infty |\log\lambda|^2 \, d\mu_{f,f}^V(\lambda) < \infty$, the Dominated Convergence Theorem allows us to obtain that
\[
\lim_{s \rightarrow 0^+} 
\left\|
\frac{\mathcal{L}_V^s f - f}{s} - (\log \mathcal{L}_V) f
\right\|_{L^2(\mathbb{R}^d)}^2 = 
\lim_{s \rightarrow 0^+}
\int_0^\infty \left| \frac{\lambda^{s}-1}{s} - \log \lambda\right|^2 \, d\mu_{f,f}^V(\lambda) = 0.
\]

Thus, we have proved that $f \in D(\LLog \mathcal{L}_V)$ and $(\LLog \mathcal{L}_V)f = (\log \mathcal{L}_V) f$.

If $f\in D(\LLog \mathcal{L}_V)$, by proceeding as in the proof of Proposition~\ref{propo: 2.2} we can see that $f\in D(\mathcal{L}_V^{s_0})\cap D(\log \mathcal{L}_V)$, for some $s_0\in (0,1)$, and $(\log \mathcal{L}_V) f = (\LLog \mathcal{L}_V) f$.

To prove part~\ref{item: teo 1.1 - b} we consider the function $F$ defined as follows:
\begin{equation*}%\label{eq: def F}
    \begin{aligned}
F: (0,\infty) & \rightarrow L^2(\mathbb{R}^d) \\
t & \mapsto \frac{e^{-t}I - T_t^V}{t} f, \text{ where } f\in L^2(\mathbb R^d). 
\end{aligned}
\end{equation*}

Since $L^2(\mathbb{R}^d)$ is a separable space, according to~\cite[p.~131]{Yo}, to show that $F$ is $L^2(\mathbb{R}^d)$-strongly measurable on $((0,\infty), dt)$ it suffices to prove that $F$ is $L^2(\mathbb{R}^d)$-weakly measurable.

Let $g\in L^2(\mathbb{R}^d)$. We define, for $t>0$,
\begin{equation*}
    H(t) = \int_{\RR^d} [F(t)](x) g(x) dx = \frac{1}{t} \left( e^{-t} \int_{\RR^d} f(x)g(x)dx - \int_{\RR^d} T_t^V(f)(x)g(x)dx\right).
\end{equation*}
Since $\{T_t^V\}_{t>0}$ is a $C_0$-semigroup of operators in $L^2(\mathbb{R}^d)$, the function $H$ is continuous in $(0,\infty)$. Hence, $H$ is a measurable function on $((0,\infty), dt)$. Thus, we proved that $F$ is weakly measurable in $((0,\infty),dt)$.


Let $m \in \mathbb{N}$. We have that 
\[
\int_{1/m}^m \|F(t)\|_{L^2(\mathbb{R}^d)} \, dt < \infty .
\]
Indeed, we can write
\begin{equation*}
    \begin{split}
        \int_{1/m}^m \|F(t)\|_{L^2(\mathbb{R}^d)} \, dt & = \int_{1/m}^m \left( \int_0^\infty \left| \frac{e^{-t} - e^{-t\lambda}}{t} \right|^2 \, d\mu_{f,f}^V(\lambda) \right)^{1/2} dt \\
        & \leq 2 \int_{1/m}^m \frac{1}{t} \left( \int_0^\infty d\mu_{f,f}^V(\lambda) \right)^{1/2} dt \\
        & \leq 4 \log(m) \|f\|_{L^2(\mathbb{R}^d)}.
    \end{split}
\end{equation*}
Hence, $F$ is $L^2(\mathbb{R}^d)$-Bochner integrable on $(1/m, m)$ for each $m\in\NN$.

Suppose now that $f \in D(\log \mathcal{L}_V)$. We are going to show that 
\begin{equation*}
    \lim_{m \rightarrow \infty} \int_{1/m}^m \frac{e^{-t}I - T_t^V}{t} f \, dt = (\log \mathcal{L}_V) f,
\end{equation*}
where the integrals are understood in the $L^2(\mathbb{R}^d)$-Bochner sense and the limit is taken in $L^2(\mathbb{R}^d)$.

Let $m \in \mathbb{N}$ and $g \in L^2(\mathbb{R}^d)$. It follows that
\begin{equation*}
    \begin{split}
        \int_{\mathbb{R}^d} |g(x)| & \int_{1/m}^m \left|
        \left( \int_0^\infty \frac{e^{-t} - e^{-t\lambda}}{t} \, dE_V(\lambda)\right) f(x) \right| dt \, dx \\
        & \leq \|g\|_{L^2(\mathbb{R}^d)} \int_{1/m}^m \left\| \left( \int_0^\infty \frac{e^{-t} - e^{-t\lambda}}{t} \, dE_V(\lambda)\right) f \right\|_{L^2(\mathbb{R}^d)} dt \\
        & \leq \|g\|_{L^2(\mathbb{R}^d)} \int_{1/m}^m 
        \left( \int_0^\infty \left|\frac{e^{-t} - e^{-t\lambda}}{t} \right|^2 \, d\mu_{f,f}^V(\lambda) \right)^{1/2} dt \\
        & \leq 4  \log (m)\|g\|_{L^2(\mathbb{R}^d)} \|f\|_{L^2(\mathbb{R}^d)} \, .
    \end{split}
\end{equation*}

Then, Fubini's theorem leads to 
\begin{equation*}
    \begin{split}
        \int_{\mathbb{R}^d} g(x) \int_{1/m}^m & \left( \int_0^\infty \frac{e^{-t} - e^{-t\lambda}}{t} \, dE_V(\lambda)\right) f(x) \, dt \, dx \\
        & = \int_{1/m}^m \int_{\mathbb{R}^d} g(x) \left( \int_0^\infty \frac{e^{-t} - e^{-t\lambda}}{t} \, dE_V(\lambda)\right) f(x) \, dx \, dt \\
        & = \int_{1/m}^m \int_0^\infty \frac{e^{-t} - e^{-t\lambda}}{t} \, d\mu_{f,g}^V(\lambda) \, dt \\
        & = \int_0^\infty \int_{1/m}^m \frac{e^{-t} - e^{-t\lambda}}{t} \, dt \, d\mu_{f,g}^V(\lambda) \\
        & = \int_{\mathbb{R}^d} g(x) \left( \int_0^\infty \int_{1/m}^m \frac{e^{-t} - e^{-t\lambda}}{t} \, dt \, dE_V(\lambda) \right) f(x) \, dx.
    \end{split}
\end{equation*}

The arbitrariness of $g$ leads to 
\begin{equation*}
    \int_{1/m}^m \left( \int_0^\infty \frac{e^{-t} - e^{-t\lambda}}{t} \, dE_V(\lambda) \right) f \, dt = \int_0^\infty \int_{1/m}^m \frac{e^{-t} - e^{-t\lambda}}{t} \, dt \, dE_V(\lambda) f.
\end{equation*}

Our objective is to show that 
\begin{equation*}
    \lim_{m \rightarrow \infty} \int_0^\infty \int_{1/m}^m \frac{e^{-t} - e^{-t\lambda}}{t} \, dt \, dE_V(\lambda) f = (\log \mathcal{L}_V) f,
\end{equation*}
where the limit is taken in $L^2(\mathbb{R}^d)$.


We have that 
\begin{equation*}
    \begin{split}
        \left\| \int_0^\infty \int_m^\infty \frac{e^{-t} - e^{-t\lambda}}{t} \, dt \, dE_V(\lambda) f \right\|_{L^2(\mathbb{R}^d)}^2 
        & = \int_0^\infty \left( \int_m^\infty \left| \frac{e^{-t} - e^{-t\lambda}}{t} \right| \, dt \right)^2 d\mu_{f,f}^V(\lambda) \\
        & \leq \int_0^1 \left( \int_m^\infty \left| \frac{e^{-t} - e^{-t\lambda}}{t} \right| \, dt \right)^2 d\mu_{f,f}^V(\lambda) \\
        & \quad + \int_1^\infty \left( \int_m^\infty \left| \frac{e^{-t} - e^{-t\lambda}}{t} \right| \, dt \right)^2 d\mu_{f,f}^V(\lambda) \\
        & =: J_1(m) + J_2(m), \quad m \in \mathbb{N}.
    \end{split}
\end{equation*}

We get, as in the proof of \cite[Theorem~2.12]{Ch-manifolds},
\begin{equation*}
    \int_m^\infty \left| \frac{e^{-t} - e^{-t\lambda}}{t} \right| dt = \int_m^\infty \frac{e^{-t\lambda} - e^{-t}}{t} \, dt \leq \int_1^\infty \frac{e^{-t\lambda} - e^{-t}}{t} \, dt \leq C(|\log\lambda| + 1),
\end{equation*}
for $0 < \lambda < 1$ and $m \in \mathbb{N}$.



Since $f \in D(\log \mathcal{L}_V)$, we have $\displaystyle \int_0^1 (|\log\lambda| + 1)^2 \, d\mu_{f,f}^V(\lambda) < \infty$, and the Dominated Convergence Theorem leads to
\[
\lim_{m \rightarrow \infty} J_1(m) = 0.
\]

On the other hand, we have that
\[
\int_m^\infty \left| \frac{e^{-t} - e^{-t\lambda}}{t} \right| \, dt \leq 2 \int_m^\infty \frac{e^{-t}}{t} \, dt \leq 2 \int_1^\infty \frac{e^{-t}}{t} \, dt, \quad \lambda > 1 \text{ and } m \in \mathbb{N}.
\]
Then, since $\mu_{f,f}^V$ is a bounded measure and $\int_1^\infty \frac{e^{-t}}{t} \, dt < \infty$, we obtain
\[
\lim_{m \rightarrow \infty} J_2(m) = 0.
\]

We conclude that 
\[ 
\lim_{m \rightarrow \infty} \left( J_1(m) + J_2(m) \right) = 0.
\]


Again, as in the proof of \cite[Theorem~2.12]{Ch-manifolds} we obtain
\[
\int_0^{1/m} \left| \frac{e^{-t} - e^{-t\lambda}}{t} \right| \, dt \leq \int_0^1 \left| \frac{e^{-t} - e^{-t\lambda}}{t} \right| \, dt \leq C(|\log \lambda| + 1), \quad \lambda > 0 \text{ and } m \in \mathbb{N}. 
\]
The Dominated Convergence Theorem leads to 
\[ 
\lim_{m \rightarrow \infty} \left\| \int_0^\infty \int_0^{1/m} \frac{e^{-t} - e^{-t\lambda}}{t} \, dt \, dE_V(\lambda) f \right\|_{L^2(\mathbb{R}^d)} = 0.
\]

We conclude that $f \in D(\Log \mathcal{L}_V)$ and 
\[ 
\lim_{m \rightarrow \infty} \int_0^\infty \int_{1/m}^m \frac{e^{-t} - e^{-t\lambda}}{t} \, dt \, dE_V(\lambda) f = \int_0^\infty \log(\lambda) \, dE_V(\lambda) f = (\log \mathcal{L}_V) f,
\]
in $L^2(\mathbb{R}^d)$, which completes the proof.

\end{proof}


We define, for every $\alpha \in (0,1)$, the operator 
\[ 
\mathcal{L}_{V,\textup{heat}}^\alpha f = \frac{1}{\Gamma(-\alpha)} \lim_{m \rightarrow \infty} \int_{1/m}^m \frac{T_t^V f - f}{t^{1+\alpha}} \, dt, \quad f \in D(\mathcal{L}_{V,\textup{heat}}^{\alpha}), 
\]
where the integrals are understood in the $L^2(\mathbb{R}^d)$-Bochner sense and the limit is taken in $L^2(\mathbb{R}^d)$. Here
\begin{equation*}
\begin{split}
    D(\mathcal{L}_{V,\textup{heat}}^{\alpha}) = & \left\{ f \in L^2(\mathbb{R}^d) : \frac{T_t^V f - f}{t^{1+\alpha}} \text{ is } L^2(\mathbb{R}^d)\text{-Bochner integrable on } ((1/m,m), dt),\right. \\
    & \left. \text{ for every } m \in \mathbb{N}, \text{ and } \lim_{m \rightarrow \infty} \int_{1/m}^m \frac{T_t^V f - f}{t^{1+\alpha}} \, dt \text{ exists in } L^2(\mathbb{R}^d) \right\}.
\end{split}
\end{equation*}


\begin{prop} \label{propo: 2.4}
    Let $f \in L^2(\mathbb{R}^d)$ and $\alpha \in (0,1)$.
    \begin{enumerate}
        \item\label{item: propo 2.4 - a} If $f \in D(\mathcal{L}_V^\alpha)$, then $f \in D(\mathcal{L}_{V,\textup{heat}}^{\alpha})$ and $\mathcal{L}_{V,\textup{heat}}^{\alpha} f = \mathcal{L}_V^\alpha f$.
        \item\label{item: propo 2.4 - b} If $f \in D(\mathcal{L}_V^\beta)$ for some $\alpha < \beta \leq 1$, then
        \[
        \mathcal{L}_V^\alpha f = \frac{1}{\Gamma(-\alpha)} \int_0^\infty \frac{T_t^V f - f}{t^{1+\alpha}} \, dt,
        \]
        where the integral converges in the $L^2(\mathbb{R}^d)$-Bochner sense.
    \end{enumerate}
\end{prop}

\begin{proof}
    We note first that part~\ref{item: propo 2.4 - a} can be proved proceeding as in the proof of part~\ref{item: teo 1.1 - b} of Theorem~\ref{thm: 1.1}. To prove part~\ref{item: propo 2.4 - b} of Proposition~\ref{propo: 2.4}, we consider the function $F_\alpha$ defined as follows:
\begin{equation*}
    \begin{aligned}
F_\alpha: (0,\infty) & \rightarrow L^2(\mathbb{R}^d) \\
t & \mapsto \frac{e^{-t}I - T_t^V}{t^{1+\alpha}} f, \text{ where } f\in L^2(\mathbb R^d). 
\end{aligned}
\end{equation*} 
    
    
    Let $f \in D(\mathcal{L}_V^\beta)$ for some $\alpha < \beta \leq 1$. We have that
    \begin{equation*}
        \begin{split}
            \int_0^\infty \|F_\alpha(t)\|_{L^2(\mathbb{R}^d)} \, dt & = \int_0^\infty \left( \int_0^\infty \left| \frac{e^{-t\lambda} - 1}{t^{1+\alpha}} \right|^2 \, d\mu_{f,f}^V(\lambda) \right)^{1/2} dt \\
            & \leq C \left( \int_0^1 \frac{dt}{t^{1+\alpha-\beta}} \left( \int_0^\infty \lambda^{2\beta} \, d\mu_{f,f}^V(\lambda) \right)^{1/2} + \int_1^\infty \frac{dt}{t^{1+\alpha}} \left(\int_0^\infty d\mu_{f,f}^V(\lambda)\right)^{1/2} \right) < \infty,
        \end{split}
    \end{equation*}
    because $\mu_{f,f}^V$ is a bounded measure and $f \in D(\mathcal{L}_V^\beta)$. 

    We conclude that $F_\alpha$ is $L^2(\mathbb{R}^d)$-Bochner integrable on $((0,\infty), dt)$. It follows that $f \in D(\mathcal{L}_{V,\textup{heat}}^\alpha)$ and that 
    \[
    \mathcal{L}_V^\alpha f = \mathcal{L}_{V,\textup{heat}}^\alpha f = \frac{1}{\Gamma(-\alpha)} \int_0^\infty \frac{T_t^V f - f}{t^{1+\alpha}} \, dt,
    \]
    which completes the proof.
\end{proof}


We define, for every $\alpha \in (0,1)$, the operator 
\[ 
\mathcal{L}_{V,\textup{heat}}^{-\alpha} f = \frac{1}{\Gamma(\alpha)} \lim_{m \rightarrow \infty} \int_{1/m}^m \frac{T_t^V f}{t^{1-\alpha}} \, dt, \quad f \in D(\mathcal{L}_{V,\textup{heat}}^{-\alpha}), 
\]
where the integrals are understood in the $L^2(\mathbb{R}^d)$-Bochner sense and the limit is taken in $L^2(\mathbb{R}^d)$. Here
\begin{equation*}
\begin{split}
    D(\mathcal{L}_{V,\textup{heat}}^{-\alpha}) = & \left\{ f \in L^2(\mathbb{R}^d) : \frac{T_t^V f}{t^{1-\alpha}} \text{ is } L^2(\mathbb{R}^d)\text{-Bochner integrable on } ((1/m,m), dt),\right. \\
    & \left. \text{ for every } m \in \mathbb{N}, \text{ and } \lim_{m \rightarrow \infty} \int_{1/m}^m \frac{T_t^V f}{t^{1-\alpha}} \, dt \text{ exists in } L^2(\mathbb{R}^d) \right\}.
\end{split}
\end{equation*}

By arguing as in Proposition~\ref{propo: 2.4} we can obtain the following result. 

\begin{prop}\label{propo: 2.5}
      Let $f \in L^2(\mathbb{R}^d)$ and $\alpha \in (0,1)$.
    \begin{enumerate}
        \item\label{item: propo 2.5 - a} If $f \in D(\mathcal{L}_V^{-\alpha})$, then $f \in D(\mathcal{L}_{V,\textup{heat}}^{-\alpha})$ and $\mathcal{L}_{V,\textup{heat}}^{-\alpha} f = \mathcal{L}_V^{-\alpha} f$.
        \item\label{item: propo 2.5 - b} If $f \in D(\mathcal{L}_V^{-\beta})$ for some $\alpha < \beta < 1$, then
        \[
        \mathcal{L}_V^{-\alpha} f = \frac{1}{\Gamma(\alpha)} \int_0^\infty \frac{T_t^V f}{t^{1-\alpha}} \, dt,
        \]
        where the integral converges in the $L^2(\mathbb{R}^d)$-Bochner sense.
    \end{enumerate}
\end{prop}

\begin{rem}
    Note that the properties established in this section can also be proved for a more general class of operators. As far as we know, these properties have not been studied carefully enough in the literature.
\end{rem}

\section{Pointwise representation of \texorpdfstring{$\log \mathcal{L}_V$}{log LV}}\label{sec: 3}

In this section we are going to prove Theorem~\ref{thm: 1.2}. 

Let $f\in C_c^\infty (\RR^d)$. According to Theorem~\ref{thm: 1.1} we have that
\[(\log \mathcal{L}_V)f=\lim_{s\to 0^+}\frac{\mathcal{L}_V^s f-f}{s}, \quad \text{in }L^2(\RR^d).\]
Then, for a certain sequence $\{s_n\}_{n\in \NN}\subseteq (0,\infty)$ such that $\lim_{n\to \infty}s_n=0$, we can write
\begin{equation}\label{eq: log LVsf ctp}
    (\log \mathcal{L}_V)f(x)=\lim_{n\to\infty}\frac{\mathcal{L}_V^{s_n} f(x)-f(x)}{s_n}, \quad \text{a.e. }x\in \RR^d.
\end{equation}
Since $f\in D(\mathcal{L}_V)$, it follows that 
\[\mathcal{L}_V^s f=\frac{1}{\Gamma(-s)}\int_0^\infty \frac{T_t^V(f)-f}{t^{s+1}}dt, \quad 0<s<1,\]
where the integral is understood in the $L^2(\RR^d)$-Bochner sense.

Indeed, let $g\in L^2(\RR^d)$. From \eqref{eq: log LVsf ctp} and the properties of the Bochner integrals, we can write
\begin{align*}
    \int_{\RR^d} g(x) \left(\int_0^\infty \frac{T_t^V(f)-f}{t^{s+1}}dt\right)(x) dx&=\int_0^\infty \int_{\RR^d} \frac{T_t^V(f)(x)-f(x)}{t^{s+1}}g(x) dx dt\\
    &=\int_0^1 \frac{T_t^V(f)(x)-f(x)}{t^{s+1}}g(x) dx dt+\int_1^\infty\frac{T_t^V(f)(x)-f(x)}{t^{s+1}}g(x) dx dt.
\end{align*}
Since $\sup_{t>0}\|T_t^V\|_{L^2(\RR^d)\to L^2(\RR^d)}<\infty$, we obtain
\begin{align*}
    \int_1^\infty\frac{\left|T_t^V(f)(x)-f(x)\right|}{t^{s+1}}&|g(x)| dx dt\\
    &\leq  \left(\sup_{t>0}\|T_t^Vf\|_{L^2(\RR^d)}+\|f\|_{L^2(\RR^d)}\right)\|g\|_{L^2(\RR^d)} \int_1^\infty \frac{dt}{t^{s+1}}\\
    &\leq \frac{C}{s} \|f\|_{L^2(\RR^d)}\|g\|_{L^2(\RR^d)}<\infty.
\end{align*}

On the other hand, since $f\in D(\mathcal L_V)$, there exists $C>0$ such that 
\[\|T_t^V(f)-f\|_{L^2(\RR^d)}\leq Ct, \quad t\in (0,1).\]
Then,
\[\int_0^1\int_{\mathbb R^d} \frac{|T_t^V(f)(x)-f(x)|}{t^{s+1}}|g(x)|dx dt\leq C \|g\|_{L^2(\RR^d)}\int_0^1 \frac{1}{t^s} dt <\infty.\]
Fubini's theorem leads to
\[\int_{\RR^d} \left(\int_0^\infty \frac{T_t^V(f)-f}{t^{s+1}}dt\right)(x) dx=\int_{\RR^d} g(x) \int_0^\infty \frac{T_t^V(f)(x)-f(x)}{t^{s+1}}dt dx.\]
The arbitrariness of $g$ implies that
\[\left(\int_0^\infty \frac{T_t^V(f)-f}{t^{s+1}}dt\right)(x)=\int_0^\infty \frac{T_t^V(f)(x)-f(x)}{t^{s+1}}dt, \quad \text{a.e. }x\in \RR^d.\]

We are going to see that, for every $x\in \RR^d$ and $r>0$,
\begin{equation}\label{eq: 2.1}
    \lim_{s\to 0^+} H_s(x)=\int_{B(x,r)} (f(x)-f(y))\int_0^\infty \frac{T_t^V(x,y)}{t}dtdy-\int_{\RR^d\setminus B(x,r)}f(y)\int_0^\infty \frac{T_t^V(x,y)}{t}dtdy-K(x,r)f(x),
\end{equation}
where
\begin{align*}
    K(x,r)&=2\log(\rho(x))+\int_0^{\rho^2(x)} \int_{\RR^d} \frac{T_t^V(x,y)-T_t(x-y)}{t}dydt-\int_0^{\rho^2(x)} \int_{\RR^d\setminus B(x,r)} \frac{T_t^V(x,y)}{t}dydt\\
    &\quad +\int_{\rho^2(x)}^\infty \int_{B(x,r)} \frac{T_t^V(x,y)}{t}dydt-\gamma,
\end{align*}
and 
\[H_s(x)=\frac1s \left(\frac{1}{\Gamma(-s)}\int_0^\infty \frac{T_t^V(f)(x)-f(x)}{t^{s+1}}dt-f(x)\right), \quad s\in (0,\tfrac12).\]

Fix $x\in \RR^d$. We can write
\begin{align*}
    \int_0^\infty \frac{T_t^V(f)(x)-f(x)}{t^{s+1}}dt&=\int_0^\infty \frac{T_t^V(f)(x)-T_ t^V(1)(x)f(x)}{t^{s+1}}dt+f(x)\int_0^\infty \frac{T_t^V(1)(x)-1}{t^{s+1}}dt\\
    &:=J_1(s)+J_2(s), \quad s\in (0,\tfrac12).
\end{align*}

According to \eqref{eq: 1.2} we have that
\begin{align*}
    \int_0^\infty \int_{\RR^d} \frac{T_t^V(x,y)}{t^{s+1}} |f(x)-f(y)| dy dt &\leq C \int_{\RR^d} |f(x)-f(y)|\int_0^\infty \frac{e^{-c|x-y|^2/t}}{t^{\frac d2 +s+1}}dt dy\\
    &\leq C\int_{\RR^d} \frac{|f(x)-f(y)|}{|x-y|^{d+2s}}dy\\
    &\leq C \left(\int_{B(x,1)} \frac{dy}{|x-y|^{d+2s-1}}+\int_{B^c(x,1)}\frac{dy}{|x-y|^{d+2s}}\right)\\
    &\leq C\left(\int_0^1 \frac{du}{u^{2s}}+\int_1^\infty \frac{du}{u^{2s+1}}\right)<\infty, \quad s\in (0,\tfrac12).
\end{align*}

Let $r>0$. It follows that
\begin{align*}
    J_1(s)&=\int_{\RR^d} (f(y)-f(x))\int_0^\infty \frac{T_t^V(x,y)}{t^{s+1}}dtdy\\
    &=\int_{B(x,r)} (f(y)-f(x))\int_0^\infty \frac{T_t^V(x,y)}{t^{s+1}}dtdy+\int_{\RR^d\setminus B(x,r)} f(y)\int_0^\infty \frac{T_t^V(x,y)}{t^{s+1}}dtdy\\
    &\quad -f(x)\int_{\RR^d\setminus B(x,r)}\int_0^\infty \frac{T_t^V(x,y)}{t^{s+1}}dtdy\\
    &:=J_{1,1}(s)+J_{1,2}(s)+J_{1,3}(s), \quad s\in (0,\tfrac12).
\end{align*}

Since $s\Gamma(-s)=-\Gamma(1-s)$, for $s\in(0,1)$, $\Gamma(-s)\sim -\frac1s$ as $s\to 0^+$. We will use this property several times in the sequel.

By proceeding as above, we get
\[\int_0^\infty \frac{T_t^V(x,y)}{t^{s+1}}dt\leq C \int_0^\infty \frac{e^{-c|x-y|^2/t}}{t^{\frac d2 +s+1}}dt \leq \frac{C}{|x-y|^{d+2s}}, \quad y\neq x.\]

If $|x-y|<r$, then $|x-y|^{d+2s}=r^{d+2s}\left|\frac{x-y}{r}\right|^{d+2s}\geq C|x-y|^{d+\frac12}$ for $s\in (0,1/4)$. %\red{El $r$ estaría fijo? Porque acá $C=C_r$; no es independiente de $r$}
Thus, since $f\in C_c^\infty(\RR^d)$, we obtain
\[|f(x)-f(y)|\int_0^\infty\frac{T_t^V(x,y)}{t^{s+1}}dt\leq \frac{C}{|x-y|^{d-\frac12}}, \quad 0<|x-y|<r.\]

Since $\int_{B(x,r)} |x-y|^{-d+\frac12}dy<\infty$, the Dominated Convergence Theorem implies that
\[\lim_{s\to 0^+} J_{1,1}(s)=\int_{B(x,r)} (f(y)-f(x)) \lim_{s\to 0^+} \int_0^\infty \frac{T_t^V(x,y)}{t^{s+1}}dt dy.\]
On the other hand, using again \eqref{eq: 1.2}, we get for every $x\neq y$ and $s\in (0,1)$ that
\[\frac{T_t^V(x,y)}{t^{s+1}}\leq C\begin{dcases*} \frac{e^{-c|x-y|^2/t}}{t^{\frac d2 +\frac32}} & $t\in (0,1)$\\ 
\frac{1}{t^{\frac d2 +1}} & $t\in (1,\infty).$
\end{dcases*}\]
Therefore, the Dominated Convergence Theorem leads to
\[\lim_{s\to 0^+} \int_0^\infty \frac{T_t^V(x,y)}{t^{s+1}}dt=\int_0^\infty \frac{T_t^V(x,y)}{t}dt,\]
so we conclude that
\[\lim_{s\to 0^+} \frac{1}{\Gamma(-s)s} J_{1,1}(s)=-\int_{B(x,r)} (f(y)-f(x))\int_0^\infty \frac{T_t^V(x,y)}{t}dt.\]

Since $\supp(f)$ is compact, by proceeding as before we obtain that
\[\lim_{s\to 0^+} \frac{1}{\Gamma(-s)s} J_{1,2}(s)=-\int_{\RR^d\setminus B(x,r)} f(y)\int_0^\infty \frac{T_t^V(x,y)}{t}dt.\]

We postpone, for the moment, the study of $J_{1,3}$ and proceed to estimate $J_2$.

We decompose $J_2$ as follows,
\begin{align*}
    J_2(s)&=f(x)\int_0^\infty \frac{T_t^V(1)(x)-T_t(1)(x)}{t^{s+1}}dt\\
    &=f(x)\int_0^{\rho^2(x)} \int_{\RR^d} \frac{T_t^V(x,y)-T_t(x-y)}{t^{s+1}}dydt +f(x)\int_{\rho^2(x)}^\infty \int_{\RR^d\setminus B(x,r)} \frac{T_t^V(x,y)}{t^{s+1}}dydt\\
    &\quad +{f(x)}\int_{\rho^2(x)}^\infty \int_{B(x,r)} \frac{T_t^V(x,y)}{t^{s+1}}dydt-f(x)\int_{\rho^2(x)}^\infty \int_{\RR^d} \frac{T_t(x-y)}{t^{s+1}}dydt\\
    &:=J_{2,1}(s)+J_{2,2}(s)+J_{2,3}(s)+J_{2,4}(s), 
\end{align*}
for $s\in \left(0,\frac{\delta}{4}\right)$, where $\delta>0$ is as in \eqref{eq: 1.5}.

Clearly, 
\[J_{2,4}(s)=-\frac{\rho(x)^{-2s}}{s}f(x), \quad s\in \left(0,\tfrac12\right).\]
Then,
\begin{align*}
    \lim_{s\to 0^+} \frac1s\left(\frac{1}{\Gamma(-s)}J_{2,4}(s)-f(x)\right)&=f(x)\lim_{s\to 0^+} \left(\frac{\rho(x)^{-2s}-1}{s}+\frac{1-\Gamma(1-s)}{s\Gamma(1-s)}\rho(x)^{-2s}\right)\\
    &=f(x)\left(-2\log(\rho(x))+\Gamma'(1)\right)=f(x)\left(-2\log(\rho(x))-\gamma\right).
\end{align*}

To estimate $J_{2,1}$, we consider $0<s<\frac{\delta}{4}$. According to \eqref{eq: 1.5}, for any $0<t<\rho^2(x)$ and $y\in \RR^d$, we can obtain
\begin{align*}
    \frac{|T_t^V(x,y)-T_t(x-y)|}{t^{s+1}}&\leq  \left(\frac{\sqrt{t}}{\rho(x)}\right)^\delta \frac{\omega_t(x-y)}{t^{s+1}} =  \frac{\omega_t(x-y)}{\left(\frac{t}{\rho^2(x)}\right)^{s+1-\delta/2}(\rho(x))^{2s+2}}\\
    &\leq \frac{\omega_t(x-y)}{\left(\frac{t}{\rho^2(x)}\right)^{1-\delta/4}(\rho(x))^{2s+2}}\\
    &= (\rho(x))^{-2s-\delta/2} \frac{\omega_t(x-y)}{t^{1-\delta/4}}\\
    &\leq \max\left\{(\rho(x))^{-\delta}, (\rho(x))^{-\delta/2}\right\} \frac{\omega_t(x-y)}{t^{1-\delta/4}}.
\end{align*}
Since $\int_0^{\rho^2(x)}\int_{\RR^d} \frac{\omega_t(x-y)}{t^{1-\delta/4}} dydt<\infty$, it follows that
\[\lim_{s\to 0^+} \frac{1}{s\Gamma(-s)}J_{2,1}(s)=-f(x)\int_0^{\rho^2(x)}\int_{\RR^d} \frac{T_t^V(x,y)-T_t(x-y)}{t}dydt.\]

We are going to combine now the estimates for $J_{1,3}$ and $J_{2,2}$. By using again \eqref{eq: 1.2}, if $0<s<\frac{\delta}{4}$, $0<t<\rho^2(x)$ and $|x-y|>r$, we get
\begin{align*}
    \frac{T_t^V(x,y)}{t^{s+1}}&\leq C \frac{e^{-c|x-y|^2/t}}{t^{d/2+s+1}}\leq C\frac{e^{-\frac{c}{2t}r^2} e^{-\frac{c}{2t}|x-y|^2}}{\left(\frac{t}{\rho^2(x)}\right)^{d/2+s+1}\rho(x)^{d+2s+2}}\\
    &\leq C \max\left\{(\rho(x))^{-d-2},(\rho(x))^{-d-2-\delta/2}\right\} \frac{e^{-\frac{c}{2t}r^2} e^{-\frac{c}{2t}|x-y|^2}}{t^{d/2+\delta/4+1}}.
\end{align*}
Noting that
\[\int_0^{\rho^2(x)} \frac{e^{-\frac{c}{2t}r^2}}{t^{d/2+\delta/4+1}}\int_{\RR^d} e^{-\frac{c}{2t}|x-y|^2} dydt\leq C \int_0^{\rho^2(x)} \frac{e^{-\frac{c}{2t}r^2}{t^{d/2+s+1}}}{t^{\delta/4+1}}dt<\infty,\]
it yields
\[\lim_{s\to 0^+} \frac{1}{s\Gamma(-s)}\left(J_{1,3}(s)+J_{2,2}(s)\right)=f(x)\int_0^{\rho^2(x)}\int_{\RR^d \setminus B(x,r)} \frac{T_t^V(x,y)}{t}dy dt.\]

Finally, for $t>\rho^2(x)$ and $0<s<\delta/4$, \eqref{eq: 1.2} leads to
\[\frac{T_t^V(x,y)}{t^{s+1}}\leq C \frac{e^{-c|x-y|^2/t}}{t^{d/2+s+1}}\leq \frac{C}{\left(\frac{t}{\rho^2(x)}\right)^{d/2+s+1}\rho(x)^{d+2s+2}}\leq \frac{C}{t^{d/2+1}}.\]
Since $\int_{\rho^2(x)}^\infty \int_{B(x,r)} t^{-d/2-1}dt<\infty$, we have
\[\lim_{s\to 0^+} \frac{1}{s\Gamma(-s)} J_{2,3}(s)=-f(x)\int_{\rho^2(x)}^\infty \int_{B(x,r)}\frac{T_t^V(x,y)}{t} dydt.\]

By combining all of the above estimates, we conclude that \eqref{eq: teo 1.2} holds. By taking $r=1$ in \eqref{eq: teo 1.2} we have established part \ref{item: teo 1.2 - a} of Theorem~\ref{thm: 1.2}.

Let us now prove \ref{item: teo 1.2 - b}. We choose $r>1$ such that $\supp(f)\subseteq B(0,r)$.

First, assume that $|x|\geq 2r$, so $f(x)=0$. Then,
\[\int_0^\infty \frac{T_t^Vf(x)-f(x)}{t^{s+1}}dt=\int_{B(0,r)}f(y)\int_0^\infty \frac{T_t^V(x,y)}{t^{s+1}}dt dy.\]
Note that for every $y\in B(0,r)$, $|x-y|\geq |x|-|y|\geq \frac{|x|}{2}\geq r>1$. Then, by using \eqref{eq: 1.2}, we get
\begin{align*}
    \left|\int_0^\infty \frac{T_t^Vf(x)-f(x)}{t^{s+1}}dt\right|&\leq C\int_{B(0,r)}|f(y)|\int_0^\infty \frac{e^{-c\frac{|x-y|^2}{t}}}{t^{\frac{d}{2}+s+1}}dtdy\\
    &\leq C \int_{B(0,r)} \frac{|f(y)|}{|x-y|^{d+2s}}dy\\
    &\leq C \frac{\|f\|_{L^1(\RR^d)}}{|x|^d}, 
\end{align*}
for every $s\in (0,1)$. It follows that, for $1<p<\infty$
\begin{equation}\label{eq: bound Lp out B}
    \frac{1}{s|\Gamma(-s)|}\left\|\int_0^\infty\frac{T_t^Vf(\cdot)-f(\cdot)}{t^{s+1}}dt\right\|_{L^p(\RR^d\setminus B(0,2r))}\leq C \|f\|_{L^1(\RR^d)} r^{-\frac{d}{p'}},
\end{equation}
and
\begin{equation}\label{eq: bound Linfty out B}
    \frac{1}{s|\Gamma(-s)|}\left\|\int_0^\infty\frac{T_t^Vf(\cdot)-f(\cdot)}{t^{s+1}}dt\right\|_{L^\infty(\RR^d\setminus B(0,2r))}\leq C \|f\|_{L^1(\RR^d)} r^{-d},
\end{equation}
for every $s\in (0,1)$. Here, the constants are independent of $s$ since $s|\Gamma(-s)|\sim 1$ as $s\to 0^+$.

Suppose now $|x|<2r$. If $|x-y|\geq  3r$, then $|y|\geq |x-y|-|x|>r$. Therefore, $\supp(f) \subseteq B(x,3r)$ and we have
\begin{align*}
    \int_{\RR^d}(f(y)-f(x))\int_0^\infty \frac{T_t^V(x,y)}{t^{s+1}}dtdy&=\int_{B(x,3r)} (f(y)-f(x))\int_0^\infty \frac{T_t^V(x,y)}{t^{s+1}}dtdy\\
    &\quad -f(x) \int_{\RR^d\setminus B(x,3r)} \int_0^\infty \frac{T_t^V(x,y)}{t^{s+1}}dtdy\\
    &:=J_{1,1}(s,x)+J_{1,3}(s,x), \quad s\in (0,1).
\end{align*}
We will prove that
\begin{equation}\label{eq: lim J11}
    \frac{J_{1,1}(s,x)}{s\Gamma(-s)}\underset{s\to 0^+}{\longrightarrow} -\int_{B(x,3r)}(f(y)-f(x))\int_0^\infty \frac{T_t^V(x,y)}{t}dtdy
\end{equation}
uniformly on $x\in B(0,2r)$. Indeed, by using \eqref{eq: 1.2}, we get
\begin{align*}
    &\left| \frac{J_{1,1}(s,x)}{s\Gamma(-s)}+\int_{B(x,3r)}(f(y)-f(x))\int_0^\infty \frac{T_t^V(x,y)}{t}dtdy\right|\\
    &\leq \int_{B(x,3r)}|f(y)-f(x)| \int_0^\infty T_t^V(x,y)\left|\frac{1}{t^{s+1}}-\frac1t\right|dtdy\\
    &\leq C \int_{B(x,3r)}|x-y| \int_0^\infty \frac{e^{-c\frac{|x-y|^2}{t}}}{t^{\frac d2}}\left|\frac{1}{t^{s+1}}-\frac1t\right|dtdy\\
    &\leq C \int_0^{3r} u^d \int_0^\infty \frac{e^{-c\frac{u^2}{t}}}{t^{\frac d2}}\left|\frac{1}{t^{s+1}}-\frac1t\right|dtdu.
\end{align*}
We observe that, for every $t\in (0,\infty)$ and $s\in (0,\tfrac12)$,
\[\frac{e^{-c\frac{u^2}{t}}}{t^{\frac d2+1}}\left|\frac{1}{t^{s}}-1\right|\leq C\left(\frac{e^{-c\frac{u^2}{t}}}{t^{\frac d2+1}}\chi_{[1,\infty)}(t)+\frac{e^{-c\frac{u^2}{t}}}{t^{\frac d2+\frac32}}\chi_{(0,1)}(t)\right).\]
Hence, we have
\begin{align*}
    \int_0^{3r} u^d \int_0^\infty \frac{e^{-c\frac{u^2}{t}}}{t^{\frac d2}}\left|\frac{1}{t^{s+1}}-\frac1t\right|dtdu &\lesssim \int_0^{3r} u^d \int_0^1 \frac{e^{-c\frac{u^2}{t}}}{t^{\frac d2+1}} dtdu+\int_0^{3r} u^d \int_1^\infty \frac{e^{-c\frac{u^2}{t}}}{t^{\frac d2+\frac32}} dtdu\\
    &\lesssim \int_0^{3r} u^d \int_0^1 \left(\frac{t}{u^2}\right)^{\frac d2+\frac14} \frac{dt}{t^{\frac d2+1}}du+\int_0^{3r} u^d \int_1^\infty \frac{1}{t^{\frac d2+\frac32}} dtdu\\
    &\lesssim \left(\int_0^{3r} \frac{du}{u^{\frac12}} \right)\left(\int_0^1 \frac{dt}{t^{\frac34}}\right)+\int_0^{3r} u^d \int_1^\infty \frac{dt}{t^{\frac32}} du<\infty.
\end{align*}
The Dominated Convergence Theorem leads to
\[\lim_{s\to 0^+} \int_0^{3r} u^d \int_0^\infty \frac{e^{-c\frac{u^2}{t}}}{t^{\frac d2}}\left|\frac{1}{t^{s+1}}-\frac1t\right|dtdu=0.\]
Thus, \eqref{eq: lim J11} is proved.

As before, we proceed to estimate $J_2(s,x)$ given below, and leave the study of $J_{1,3}(s,x)$ for later. We write
\begin{align*}
    J_2(s,x)&=f(x)\int_0^\infty \frac{T_t^V(1)(x)-T_t(1)(x)}{t^{s+1}} dt\\
    &=f(x)\int_0^{\rho^2(x)} \int_{\RR^d} \frac{T_t^V(x,y)-T_t(x-y)}{t^{s+1}}dydt +f(x)\int_{\rho^2(x)}^\infty \int_{\RR^d\setminus B(x,r)} \frac{T_t^V(x,y)}{t^{s+1}}dydt\\
    &\quad +f(x)\int_{\rho^2(x)}^\infty \int_{B(x,r)} \frac{T_t^V(x,y)}{t^{s+1}}dydt-f(x)\int_{\rho^2(x)}^\infty \int_{\RR^d} \frac{T_t(x-y)}{t^{s+1}}dydt\\
    &:=J_{2,1}(s,x)+J_{2,2}(s,x)+J_{2,3}(s,x)+J_{2,4}(s,x)
\end{align*}
for $s\in (0,\frac{\delta}{4})$.

By \eqref{eq: 1.5} we get
\begin{align*}
    \left|\frac{J_{2,1}(s,x)}{s\Gamma(-s)}+\int_0^{\rho^2(x)} \int_{\RR^d} \frac{T_t^V(x,y)-T_t(x-y)}{t}dydt\right|&\leq \int_0^{\rho^2(x)} \int_{\RR^d} |T_t^V(x,y)-T_t(x-y)|\left|\frac{1}{t^{s+1}}-\frac1t\right|dydt\\
    &\lesssim \int_0^{\rho^2(x)}\left|\frac{1}{t^{s+1}}-\frac1t\right| \int_{\RR^d} \left(\frac{\sqrt{t}}{\rho(x)}\right)^\delta \omega_t(x-y) dydt\\
    &\lesssim \int_0^{\rho^2(x)} \left(\frac{\sqrt{t}}{\rho(x)}\right)^\delta \left|\frac{1}{t^{s}}-1\right|\frac{dt}{t},
\end{align*}
for every $s\in (0,\frac{\delta}{4})$.

Suppose $\mathbb K$ is some compact subset of $\RR^d$. Then, there exist $x_1,\dots, x_m\in \mathbb K$ such that $\mathbb K\subseteq \bigcup_{j=1}^m B(x_j,\rho(x_j))$. Therefore, from \eqref{eq: rox_vs_roy}, for every $y\in B(x_j,\rho(x_j))$ and $j=1, \dots, m$, we can obtain that there exist constants $0<K_1<K_2<\infty$ for which $K_1\leq \rho(y)\leq K_2$ for every $y\in \mathbb K$.

Applying this property on $\overline{B(0,2r)}$, we get that, for every $s\in (0,\frac{\delta}{4})$,
\begin{equation*}
    \left|\frac{J_{2,1}(s,x)}{s\Gamma(-s)}+\int_0^{\rho^2(x)} \int_{\RR^d} \frac{T_t^V(x,y)-T_t(x-y)}{t}dydt\right|\lesssim \frac{1}{K_1^\delta} \int_0^{A} t^{\frac{\delta}{2}-1}\left|\frac{1}{t^{s}}-1\right|dt
\end{equation*}
for $A=K_2^2>0$.

Since 
\begin{align}\label{eq: integral 0 a A}
    \int_0^{A} t^{\frac{\delta}{2}-1}\left|\frac{1}{t^{s}}-1\right|dt& \leq \int_0^{A} t^{\frac{\delta}{2}-1}\left(\frac{1}{t^{s}}+1\right)dt\leq \int_0^{A} t^{\frac{\delta}{2}-1}+\frac{1}{A^s}\int_0^{A} \frac{t^{\frac{\delta}{2}-1}}{\left(\frac{t}{A}\right)^{s}}dt\nonumber\\
    &\leq \int_0^{A} t^{\frac{\delta}{2}-1}+A^{\frac{\delta}{4}-s}\int_0^{A} t^{\frac{\delta}{4}-1} dt <\infty
\end{align}
for every $s\in (0,\frac{\delta}{4})$, we conclude that
\begin{equation}\label{eq: lim J21}
    \frac{J_{2,1}(s,x)}{s\Gamma(-s)}\underset{s\to 0^+}{\longrightarrow} -\int_0^{\rho^2(x)} \int_{\RR^d} \frac{T_t^V(x,y)-T_t(x-y)}{t}dydt
\end{equation}
uniformly in $x\in B(0,2r)$.

We also have that 
\begin{align*}
    \left|\frac{J_{1,3}(s,x)+J_{2,2}(s,x)}{s\Gamma(-s)}+\int_0^{\rho^2(x)}\int_{\RR^d\setminus B(x,3r)}\frac{T_t^V(x,y)}{t}dydt\right|&\lesssim \int_0^{\rho^2(x)}\int_{\RR^d\setminus B(x,3r)}T_t^V(x,y)\left|\frac{1}{t^{s+1}}-\frac1t\right|dydt\\
    &\lesssim \int_0^{\rho^2(x)}\int_{\RR^d\setminus B(x,3r)}\frac{e^{-c\frac{|x-y|^2}{t}}}{t^{\frac d2}}\left|\frac{1}{t^{s+1}}-\frac1t\right|dydt\\
    &\lesssim \int_0^{\rho^2(x)}\frac{e^{-c\frac{r^2}{t}}}{t}\left|\frac{1}{t^{s}}-1\right|dt\\
    &\lesssim r^{-\delta}\int_0^A t^{\frac{\delta}{2}-1}\left|\frac{1}{t^{s}}-1\right|dt<\infty
\end{align*}
proceeding as in \eqref{eq: integral 0 a A}. Therefore,
\begin{equation}\label{eq: lim J13+J22}
    \frac{J_{1,3}(s,x)+J_{2,2}(s,x)}{s\Gamma(-s)}\underset{s\to 0^+}{\longrightarrow}-\int_0^{\rho^2(x)}\int_{\RR^d\setminus B(x,3r)}\frac{T_t^V(x,y)}{t}dydt
\end{equation}
uniformly in $x\in B(0,2r)$.

To estimate the limit for $J_{2,3}(s,x)$ as $s\to 0^+$, we can write
\begin{align*}
    \left|\frac{1}{s\Gamma(-s)} J_{2,3}(s,x)+\int_{\rho^2(x)}^\infty\int_{B(x,3r)}\frac{T_t^V(x,y)}{t}dydt\right|&\lesssim \int_0^{\rho^2(x)}\int_{B(x,3r)}\frac{e^{-c\frac{|x-y|^2}{t}}}{t^{\frac d2}}\left|\frac{1}{t^{s+1}}-\frac1t\right|dydt\\
    &\lesssim r^d\int_A^\infty \frac{1}{t^{\frac d2+1}}\left|\frac{1}{t^s}-1\right|dt, 
\end{align*}
where, as before, the constants involved do not depend on $x$. Since
\[ \frac{1}{t^{\frac d2+1}}\left(\frac{1}{t^s}+1\right)= \frac{1}{t^{\frac d2+1}}\left(\frac{1}{A^s\left(\frac{t}{A}\right)^s}+1\right)\leq C(A) \frac{1}{t^{\frac d2+1}}\]
for every $s\in (0,1)$ and $t>A$. Hence,
\[\lim_{s\to 0^+} \int_A^\infty \frac{1}{t^{\frac d2+1}}\left|\frac{1}{t^s}-1\right|dt=0\]
and
\begin{equation}\label{eq: lim J23}
    \frac{1}{s\Gamma(-s)} J_{2,3}(s,x)\underset{s\to 0^+}{\longrightarrow} -\int_{\rho^2(x)}^\infty\int_{B(x,3r)}\frac{T_t^V(x,y)}{t}dydt
\end{equation}
uniformly in $x\in B(0,2r)$.

Finally, it is left to prove
\begin{equation}\label{eq: lim J24}
    \frac{\frac{1}{\Gamma(-s)} J_{2,4}(s,x)-1}{s}\underset{s\to 0^+}{\longrightarrow} -2\log(\rho(x))-\gamma
\end{equation}
uniformly in $x\in B(0,2r)$. 

By using the differential mean value theorem, we have
\[\frac{u^{-2s}-1}{u}=-2\log (u) u^{-2\xi}\]
for some $\xi\in (0,s)$, with $s\in (0,1)$ and $u>0$. Let $0<a<b<\infty$. We have that
\[b^{-2z}-1\leq u^{-2z}-1\leq a^{-2z}-1, \quad z>0, \ a\leq u\leq b.\]
Hence, $u^{-2z}\to 1$ as $z\to 0^+$ uniformly in $u\in[a,b]$. This yields $\frac{u^{-2s}-1}{u}\to -2\log(u)$ as $s\to 0^+$, uniformly in $u\in[a,b]$. 

As it was proved above, since $K_1\leq \rho(x)\leq K_2$ for every $x\in B(0,2r)$, we conclude that \eqref{eq: lim J24} holds.

By combining \eqref{eq: lim J11}, \eqref{eq: lim J21}, \eqref{eq: lim J13+J22}, \eqref{eq: lim J23} and \eqref{eq: lim J24}, we obtain that
\begin{align}\label{eq: lim uniform}
    \frac1s \left(\frac{1}{\Gamma(-s)}\int_0^\infty \frac{T_t^V(f)(x)-f(x)}{t^{s+1}}dt-1\right)\underset{s\to 0^+}{\longrightarrow}& \int_{B(x,r)}(f(x)-f(y))\int_0^\infty \frac{T_t^V(x,y)}{t}dtdy\nonumber\\
    & -\int_{\RR^d\setminus B(x,3r)} f(y)\int_0^\infty \frac{T_t^V(x,y)}{t}dtdy-K(x,3r)f(x)
\end{align}
uniformly in $x\in B(0,2r)$.

Let us prove the convergence on $L^p(\RR^d)$ for $1<p\leq \infty$. Let $\epsilon>0$. According to \eqref{eq: bound Lp out B} and \eqref{eq: bound Linfty out B}, there exists $r_0>0$ such that
\[\frac{1}{s|\Gamma(-s)|}\left\|\int_0^\infty \frac{T_t^V(f)(\cdot)-f(\cdot)}{t^{s+1}}dt\right\|_{L^p(\RR^d\setminus B(0,2r_0))}<\epsilon, \quad s\in (0,1).\]

By \eqref{eq: lim uniform}, there exists $s_0\in (0,1)$ for which
\begin{align*}
    &\left\|\frac1s \left(\frac{1}{\Gamma(-s)}\int_0^\infty \frac{T_t^V(f)(\cdot)-f(\cdot)}{t^{s+1}}dt-1\right)-\int_{B(\cdot,r)}(f(\cdot)-f(y))\int_0^\infty \frac{T_t^V(\cdot,y)}{t}dtdy\right. \\
    & \quad \left.-\int_{\RR^d\setminus B(\cdot,3r)} f(y)\int_0^\infty \frac{T_t^V(\cdot,y)}{t}dtdy-K(\cdot,3r)f(\cdot)\right\|_{L^p(B(0,2r_0))}<\epsilon, \quad 0<s<s_0.
\end{align*}
Therefore,
\[\lim_{s\to 0^+} \frac1s \left(\frac{1}{\Gamma(-s)}\int_0^\infty \frac{T_t^V(f)-f}{t^{s+1}}dt-1\right)=(\log \mathcal{L}_V)f, \quad \text{in }L^p(\RR^d),\]
and Theorem~\ref{thm: 1.2} is now proved.

The results in Theorem~\ref{thm: 1.2} are also valid when we take the limit on $s$ to zero from the left.


\begin{cor}\label{coro: 2.1}
    Let $d\geq 3$ and $V\in \RH_q$ with $q>\frac d2$. If $f\in \Lip^\theta (\RR^d)$ for some $\theta\in (0,1]$ and $\int_{\RR^d} |f(y)|(1+|y|)^{-d}dy<\infty$, then
    \[(\log \mathcal{L}_V)(f)(x)=-\lim_{h\to 0^+}\frac{\mathcal{L}^{-h}_V-I}{h}(f)(x), \quad  x\in \RR^d.\]
\end{cor}

\begin{proof}Let $h>0$ and $x\in \RR^d$. We can write
\begin{align*}
    \mathcal{L}^{-h}_Vf(x)&-f(x)\\
    %&=\frac{1}{\Gamma(h)}\int_0^{\rho^2(x)}\left(T_t^V(f)(x)-T_t^V(1)(x)f(x)\right)t^{h-1} dt\\
    %&\quad +\frac{1}{\Gamma(h)}\int_0^{\rho^2(x)}\left(T_t^V(1)(x)-T_t(1)(x)\right)f(x)t^{h-1} dt\\
    %&\quad +\frac{1}{\Gamma(h)}\int_{\rho^2(x)}^\infty T_t^V(f)(x)t^{h-1} dt+\frac{1}{\Gamma(h)}\int_0^{\rho^2(x)} T_t(1)(x)t^{h-1} dt f(x)-f(x)\\
    &=\frac{1}{\Gamma(h)}\int_0^{\rho^2(x)}\left(T_t^V(f)(x)-T_t^V(1)(x)f(x)\right)t^{h-1} dt+\frac{1}{\Gamma(h)}\int_{\rho^2(x)}^\infty T_t^V(f)(x)t^{h-1} dt\\
    &\quad +\frac{f(x)}{\Gamma(h)}\int_0^{\rho^2(x)} \left(T_t^V(1)(x)-T_t(1)(x)\right)t^{h-1}dt+\left(\frac{1}{\Gamma(h)}\int_0^{\rho^2(x)} T_t(1)(x)t^{h-1} dt-1\right)f(x)\\
    &:=A_1(x,h)+A_2(x,h)+A_3(x,h)+A_4(x,h).
\end{align*}

By using that $f\in \Lip^\theta(\RR^d)$ and (\ref{eq: 1.2}) we obtain
\begin{align*}
 T_t^V(x,y)|f(y)-f(x)|t^{h-1}&\le C|x-y|^\theta\frac{e^{-c\frac{|x-y|^2}{t}}}{t^{d/2+1-h}}\\
 &\le C\frac{e^{-\frac c2\frac{|x-y|^2}{t}}}{t^{d/2}}\left(\frac{\rho^2(x)}{t}\right)^{1-h-\theta/2}\rho(x)^{\theta-2+2h}\\
 &\le C\frac{e^{-\frac c2\frac{|x-y|^2}{t}}}{t^{d/2}}\left(\frac{\rho^2(x)}{t}\right)^{1-\theta/2}\max\{\rho^{\theta-2}(x),\rho^{\theta}(x)\}\\
 &\le C\frac{e^{-\frac c2\frac{|x-y|^2}{t}}}{t^{1+(d-\theta)/2}}\max\{1,\rho^2(x)\},
    \end{align*}
when $h\in (0,1)$, $y\in\mathbb{R}^d$ and $0<t<\rho^2(x)$.

Since 
\[
\int_0^{\rho^2(x)}\int_{\mathbb{R}^d}\frac{e^{-\frac c2\frac{|x-y|^2}{t}}}{t^{(d-\theta)/2+1}}dydt\le C\int_0^{\rho^2(x)}t^{\theta/2-1}dt\le C\rho^\theta(x),
\]
the Dominated Convergence Theorem leads to
\[
\lim_{h\to 0^+}\frac{A_1(x,h)}{h}=\int_0^{\rho^2(x)}\frac{T_t^V(f)(x)-T_t^V(1)(x)f(x)}{t}dt.
\]

%\begin{align*}
    %|A_1(x,h)|&\leq \frac{1}{\Gamma(h)}\int_0^{\rho^2(x)}\int_{\RR^d}T_t^V(x,y)|f(y)-f(x)|t^{h-1} dydt\\
    %&\leq C\frac{1}{\Gamma(h)} \int_0^{\rho^2(x)}\int_{\RR^d}|x-y|^\theta \frac{e^{-c\frac{|x-y|^2}{t}}}{t^{d/2-h+1}} dydt\\
    %&\leq C\frac{1}{\Gamma(h)} \int_0^{\rho^2(x)}\int_{\RR^d} \frac{e^{-c\frac{|x-y|^2}{2t}}}{t^{d/2}} dy \frac{dt}{t^{1-h-\theta/2}}\\
    %&\leq C\frac{\rho^{2h+\theta}(x)}{\Gamma(h)(2h+\theta)}\\
    %&\leq C\rho^{\theta}(x)
%\end{align*}
%for $h\in (0,1)$.

By using that $f\in \Lip^\theta(\mathbb{R}^d)$ and (\ref{eq: 1.3}) we get
\begin{align*}
T_t^V(x,y)|f(y)|t^{h-1}&\le {C_N}\frac{e^{-c\frac{|x-y|^2}{t}}}{t^{1-h+d/2}} \left(\frac{\rho(x)}{\rho(x)+\sqrt{t}}\right)^{{N}}(|x-y|^\theta+|f(x)|)\\
&\le {C_N}\rho(x)^{{N}}\frac{e^{-\frac c2\frac{|x-y|^2}{t}}}{t^{{N/2+1}-h+d/2}}(t^{\theta/2}+|f(x)|)\\
&\le {C_N}\rho(x)^{{N}}\frac{e^{-\frac c2\frac{|x-y|^2}{t}}}{t^{d/2}}\left(\frac{\rho^2(x)}{t}\right)^{{N/2+1}-h}(t^{\theta/2}+|f(x)|)\rho(x)^{2h-{N+2}}\\
&\le {C_N}\rho(x)^{{N}}\frac{e^{-\frac c2\frac{|x-y|^2}{t}}}{t^{d/2}}\left(\frac{\rho^2(x)}{t}\right)^{{N/2}}(t^{\theta/2}+|f(x)|)\rho(x)^{2h-{N+2}}\\
&\le {C_N}\rho(x)^{{N-2}+2h}\frac{e^{-\frac c2\frac{|x-y|^2}{t}}}{t^{d/2}}\frac{1}{t^{{N/2}}}(t^{\theta/2}+|f(x)|)\\
%&\le C max\{\rho(x)^{10},\rho(x)^8\}\frac{e^{-c\frac{|x-y|^2}{t}}}{t^{d/2}}\frac{1}{t^5}(t^{\theta/2}+|f(x)|),
\end{align*}
for $h\in (0,1)$, $y\in \mathbb{R}^d$ and $t\ge \rho^2(x)$.

{By choosing $N>3$,}
\[
\int_{\rho^2(x)}^\infty \int_{\mathbb{R}^d}\frac{e^{-\frac c2\frac{|x-y|^2}{t}}}{t^{d/2}}\frac{1}{t^{{N/2}}}(t^{\theta/2}+|f(x)|)dydt\le C\int_{\rho^2(x)}^\infty\frac{t^{\theta/2}+|f(x)|}{t^{{N/2}}}dt<\infty.
\]
According to Dominated Convergence Theorem we get
\[
\lim_{h\to 0^+}\frac{A_2(x,h)}{h}=\int_{\rho^2(x)}^\infty\frac{T_t^V(f)(x)}{t}dt.
\]

%For $A_2$, we have
%\begin{align*}
  % |A_2(x,h)|&\leq \frac{C}{\Gamma(h)}\int_{\rho^2(x)}^\infty  \int_{\RR^d}\frac{e^{-c\frac{|x-y|^2}{t}}}{t^{d/2-h+1}}(x) |f(y)| dy dt\\
  % &\leq \frac{C}{\Gamma(h)} \left(\int_{\RR^d}\frac{|f(y)|}{|x-y|^{\sigma}}dy\right)\left(\int_{\rho^2(x)}^\infty \frac{dt}{t^{d/2-h-\sigma/2+1}}\right)...
%\end{align*}

By (\ref{eq: 1.5}) it follows that
\begin{align*}
 |T_t^V(x,y)-T_t(x-y)|t^{h-1}&\le C\frac{1}{\rho(x)^\delta}\frac{\omega_t(x-y)}{t^{1-h-\delta/2}}\\
& \le C\rho(x)^{2h-2} \omega_t(x-y) \left(\frac{\rho^2(x)}{t}\right)^{1-h-\delta/2}\\
&\le C\rho(x)^{2h-\delta}\omega_t(x-y) t^{-1+\delta/2}\\
&\le C\max\{\rho(x)^{2-\delta},\rho(x)^{-\delta}\}\omega_t(x-y)t^{-1+\delta/2},
\end{align*}
when $h\in (0,1)$, $y\in \mathbb{R}^d$ and $0<t<\rho^2(x)$.

Since
\[
\int_0^{\rho^2(x)}\int_{\mathbb{R}^d}\omega_t(x-y)t^{-1+\delta/2}dydt\le C\rho(x)^\delta,
\]
the Dominated Convergence Theorem leads to
\[
\lim_{h\to 0^+}\frac{A_3(x,h)}{h}=f(x)\int_0^{\rho^2(x)}\frac{T_t^V(1)(x)-T_t(1)(x)}{t}dt.
\]


%We now estimate $A_3$. 
%\begin{align*}
    %|A_3(x,h)|&\leq \frac{C|f(x)|}{\Gamma(h)\rho^\delta(x)} \int_0^{\rho^2(x)} \int_{\RR^d} \frac{\omega_t(x-y)}{t^{1-h-\delta/2}} dy dt\\
   % &\leq \frac{C|f(x)|}{\Gamma(h)\rho^\delta(x)} \int_0^{\rho^2(x)}  \frac{dt}{t^{1-h-\delta/2}} dt\\
   % &\leq \frac{C|f(x)|}{\Gamma(h)(2h+\delta)}\\
    %&\leq C|f(x)|
%\end{align*}
%for $h\in (0,1)$.

Finally, we have that
\begin{align*}
\frac{A_4(x,h)}{h}&=\frac{f(x)}{h}\left(\frac{\rho(x)^{2h}}{h\Gamma(h)}-1\right)=f(x)\left(\frac{\rho(x)^{2h}-1}{\Gamma(h+1)h}+ \frac{1-\Gamma(h+1)}{\Gamma(h+1)h}\right), \quad h\in (0,1).
\end{align*}
Then,
\[
\lim_{h\to 0^+}\frac{A_4(x,h)}{h}=f(x)( \gamma+2\log\rho(x)).
\]


%Finally, for $A_4$,
%\begin{equation*}
    %|A_4(x,h)|\leq \frac{C|f(x)|}{\Gamma(h)} \int_0^{\rho^2(x)} t^{h-1}  dt\leq \frac{C|f(x)|\rho^{2h}(x)}{2h\Gamma(h)}\leq C|f(x)|
%\end{equation*}
%for $h\in (0,1)$.

We conclude that 
\begin{align*}
-\lim_{h\to 0^+}\frac{\mathcal{L}_V^{-h}(f)(x)-f(x)}{h}&=\int_0^{\rho^2(x)}\frac{T_t^V(1)f(x)-T_t^V(f)(x)}{t}dt\\
&-\int_{\rho^2(x)}^\infty \frac{T_t^V(f)(x)}{t}dt\\
&+f(x)\int_0^{\rho^2(x)}\frac{T_t(1)(x)-T_t^V(1)(x)}{t}dt\\
&-f(x)(\gamma+2\log\rho(x)).
\end{align*}


    By a rearrangement of the terms and taking into account that $\int_{\mathbb{R}^d}|f(y)|(1+|y|)^{-d}<\infty$,  Theorem~\ref{thm: 1.2} and the comments after this one allow us to conclude that
    \[-\lim_{h\to 0^+}\frac{\mathcal{L}^{-h}_V(f)(x)-f(x)}{h}=\lim_{h\to 0^+}\frac{\mathcal{L}^{h}_V(f)(x)-f(x)}{h}=(\log \mathcal{L}_V)(f)(x).\qedhere\]
\end{proof}

\section{Proof of Theorem~\ref{thm: 1.3}}\label{sec: 4}


We define $u(x,t) = \mathcal{L}_V^{-t} f(x)$, $x \in \mathbb{R}^d$ and $t > 0$. Since $f \in L^1(\mathbb{R}^d)$, for every $0 < t < d/2$, $\mathcal{L}_V^{-t} f \in L^{d/(d-2t),\infty}(\mathbb{R}^d)$. Moreover, since $f \in \Lip_V^\theta$, we have $\mathcal{L}_V^{-t} f \in \Lip_V^{\theta + t}$ for $t \in (0, 1-\theta)$. Then, $\mathcal{L}_V^{-t} f \in \Lip^{\theta + t}$ and $\| \mathcal{L}_V^{-t} f (\cdot) \rho(\cdot)^{-(\theta+t)} \|_\infty < \infty$ for every $t \in (0, 1-\theta)$. It follows that

\begin{equation*}
    \int_{\mathbb{R}^d} \frac{|\mathcal{L}_V^{-t} f(x)|}{(1+|x|)^d} \, dx \leq \| \mathcal{L}_V^{-t} f(\cdot) \rho(\cdot)^{-(\theta+t)} \|_{\infty}
    \int_{\mathbb{R}^d} \frac{\rho(x)^{\theta+t}}{(1+|x|)^d} \, dx, \quad t \in (0, 1-\theta).
\end{equation*}



According to the comments after Theorem~\ref{thm: 1.2} and Corollary~\ref{coro: 2.1}, we obtain
\begin{equation}\label{eq: 3.1}
    \Log (\mathcal{L}_V^{-t} f)(x) = \lim_{h \rightarrow 0^+} \frac{\mathcal{L}_V^{-h} - I}{h} \mathcal{L}_V^{-t} f(x), \quad x \in \mathbb{R}^d \text{ and } t \in (0, 1-\theta).
\end{equation}
We are going to show that
\begin{equation*}
    \mathcal{L}_V^{-h} (\mathcal{L}_V^{-t} f)(x) = \mathcal{L}_V^{-(h+t)} f(x), \quad t > 0, \, x \in \mathbb{R}^d \text{ and } 0 < h < 1 - t - \theta.
\end{equation*}

Assume first that $0 < h + t + \theta < 1$, with $h, t > 0$. We have that
\begin{equation*}
    \mathcal{L}_V^{-t} f(x) = \frac{1}{\Gamma(t)} \int_{\mathbb{R}^d} f(y) \int_0^\infty T_u^V(x,y) u^{t-1} \, du \, dy, \quad x \in \mathbb{R}^d.
\end{equation*}

We can write
\begin{equation*}
    \begin{split}
        \int_{\mathbb{R}^d} \int_{\mathbb{R}^d} |f(y)| & \int_{0}^\infty T_u^V(z,y) u^{t-1} \, du \int_0^{\infty} T_s^V(x,z) s^{h-1} \, ds \, dy \, dz \\
        & = \int_{\mathbb{R}^d} |f(y)| \int_{0}^\infty \int_0^{\infty} \int_{\mathbb{R}^d} T_u^V(z,y) T_s^V(x,z) \, dz \, u^{t-1} s^{h-1} \, du \, ds \, dy \\
        & = \int_{\mathbb{R}^d} |f(y)| \int_{0}^\infty \int_0^{\infty} T_{u+s}^V(x,y) u^{t-1} s^{h-1} \, du \, ds \, dy \\
        & = \int_{\mathbb{R}^d} |f(y)| \int_{0}^\infty \int_u^{\infty} T_{v}^V(x,y) (v-u)^{h-1} \, dv \, u^{t-1} \, du \, dy \\
        & = \int_{\mathbb{R}^d} |f(y)| \int_{0}^\infty \int_0^{v} (v-u)^{h-1} u^{t-1} \, du \, T_{v}^V(x,y) \, dv \, dy \\
        & = \frac{\Gamma(h)\Gamma(t)}{\Gamma(t+h)} \int_{\mathbb{R}^d} |f(y)| \int_{0}^\infty v^{h+t-1} T_{v}^V(x,y) \, dv \, dy \\
        & = \Gamma(h)\Gamma(t) \, \mathcal{L}_V^{-h-t}(|f|)(x) < \infty, \quad x \in \mathbb{R}^d.
    \end{split}
\end{equation*}

This justifies the interchange in the order of integration, obtaining 
\begin{equation}\label{eq: 3.2}
    \mathcal{L}_V^{-h}(\mathcal{L}_V^{-t} f)(x) = \mathcal{L}_V^{-(t+h)} f(x), \quad x \in \mathbb{R}^d.
\end{equation}

According to~\eqref{eq: 3.1} we get
\begin{equation*}
    \lim_{h \rightarrow 0^+} \frac{\mathcal{L}_V^{-(t+h)} f(x) - \mathcal{L}_V^{-t} f(x)}{h} = -\log(\mathcal{L}_V) (\mathcal{L}_V^{-t} f)(x), \quad x \in \mathbb{R}^d \text{ and } t \in (0, 1-\theta).
\end{equation*}
We are going to show that
\begin{equation}\label{eq: 3.3}
    \lim_{h \rightarrow 0^+} \frac{\mathcal{L}_V^{-(t-h)} f(x) - \mathcal{L}_V^{-t} f(x)}{-h} = -\log(\mathcal{L}_V) (\mathcal{L}_V^{-t} f)(x), \quad x \in \mathbb{R}^d \text{ and } t \in (0, 1-\theta).
\end{equation}

By using~\eqref{eq: 3.2} we get
\begin{equation*}
    \mathcal{L}_V^{-(t-h)} f(x) - \mathcal{L}_V^{-t} f(x) = \mathcal{L}_V^{-(t-h)} f(x) - \mathcal{L}_V^{-h} (\mathcal{L}_V^{-(t-h)} f)(x) = \frac{I - \mathcal{L}_V^{-h}}{h} \mathcal{L}_V^{-(t-h)} f(x),
\end{equation*}
for $x \in \mathbb{R}^d$ and $0 < h < t < 1-\theta$. 

We now prove that 
\begin{equation}\label{eq: 3.4}
\lim_{h \rightarrow 0^+} \mathcal{L}_V^{-(t-h)} f(x) = \mathcal{L}_V^{-t} f(x), \quad x \in \mathbb{R}^d \text{ and } 0 < t < 1-\theta.     
\end{equation}
Let $x \in \mathbb{R}^d$ and $0 < t < 1-\theta$. We decompose $\mathcal{L}_V^{-(t-h)} f(x)$ as follows:
\begin{equation*}
    \begin{split}
        \mathcal{L}_V^{-(t-h)} f(x) & = \frac{1}{\Gamma(t-h)} \left[
        \int_0^1 u^{t-h-1} \int_{\mathbb{R}^d} T_u^V(x,y) f(y) \, dy \, du + \int_1^\infty u^{t-h-1} \int_{\mathbb{R}^d} T_u^V(x,y) f(y) \, dy \, du \right] \\
        & =: I_1(h) + I_2(h), \quad 0 < h < t.
    \end{split}
\end{equation*}

Since $f \in \textup{Lip}_V^{\theta}$, by using~\eqref{eq: 1.2} we get, {for $0 < h < t/2$},
\begin{equation}\label{eq: 3.5}
    \begin{split}
        \int_0^1 u^{t-h-1} \int_{\mathbb{R}^d} T_u^V(x,y) |f(y)| \, dy \, du & \leq C \int_0^1 u^{t-h-1} \int_{\mathbb{R}^d} T_u(x-y) |f(y)| \, dy \, du \\
        & \leq C \int_0^1 u^{t/2-1} \int_{\mathbb{R}^d} T_u(x-y) |f(y)| \, dy \, du \\
        & \leq C \left( \int_{|x-y|<1} \frac{|f(y)|}{|x-y|^{d-t}} \, dy + \int_{\mathbb{R}^d} |f(y)| \, dy \right) \\
        & \leq C \left( \int_{|x-y|<1} \frac{|f(y)-f(x)| + |f(x)|}{|x-y|^{d-t}} \, dy + \int_{\mathbb{R}^d} |f(y)| \, dy \right) < \infty,
    \end{split}
\end{equation}
where $C$ does not depend on $h$. Then, the Dominated Convergence Theorem leads to
\begin{equation*}
    \lim_{h \rightarrow 0^+} I_1(h) = \frac{1}{\Gamma(t)} \int_0^1 u^{t-1} \int_{\mathbb{R}^d} T_u^V(x,y) f(y) \, dy \, du.
\end{equation*}

Since $f \in L^1(\mathbb{R}^d)$, by using again~\eqref{eq: 1.2} we obtain
\begin{equation}\label{eq: 3.6}
    \begin{split}
        \int_1^\infty u^{t-h-1} \int_{\mathbb{R}^d} T_u^V(x,y) |f(y)| \, dy \, du  & \leq C \int_1^\infty u^{t-h-1-d/2} \, du \int_{\mathbb{R}^d} |f(y)| \, dy \\
        & \leq C \int_1^\infty u^{t-1-d/2} \, du \int_{\mathbb{R}^d} |f(y)| \, dy < \infty, \quad 0 < h < t, 
    \end{split}
\end{equation}
where again $C$ does not depend on $h$. According to the Dominated Convergence Theorem we get
\begin{equation*}
    \lim_{h \rightarrow 0^+} I_2(h) = \frac{1}{\Gamma(t)} \int_{1}^{\infty} u^{t-1} \int_{\mathbb{R}^d} T_u^V(x,y) f(y) \, dy \, du,
\end{equation*}
and~\eqref{eq: 3.4} is proved.

By following carefully the proofs of Lemma 4.1 and Theorem 1.6 in~\cite{dLT}, we can deduce that if $0 < a < b < +\infty$ and $f \in \textup{Lip}^\theta_V$ with $\theta + b < 1$, there exists $C > 0$ such that
\begin{equation}\label{eq: 3.7}
    \|\mathcal{L}_V^{-s} f\|_{\textup{Lip}_V^{\theta+s}} \leq C \|f\|_{\textup{Lip}_V^{\theta}}, \quad s \in [a,b],
\end{equation}
where $C$ does not depend on $s \in [a,b]$. 

Let $x \in \mathbb{R}^d$ and $0 < t < 1-\theta$. By decomposing as in the proof of Theorem~\ref{thm: 1.2} and Corollary~\ref{coro: 2.1}, we consider
\begin{equation*}
    \begin{split}
        J_1(h) & = \int_{B(x,1)} \left( \mathcal{L}_V^{-(t-h)} f(x) - \mathcal{L}_V^{-(t-h)} f(y) \right) \int_0^\infty T_u^V(x,y) u^{h-1} \, du \, dy, \\
        J_2(h) & = \int_{\mathbb{R}^d \setminus B(x,1)} \mathcal{L}_V^{-(t-h)} f(y) \int_0^\infty T_u^V(x,y) u^{h-1} \, du \, dy, \\
        J_3(h) & = \left( - \int_{\mathbb{R}^d \setminus B(x,1)} \int_0^\infty T_u^V(x,y) u^{h-1} \, du \, dy 
        + \int_{0}^{\rho^2(x)} \int_{\mathbb{R}^d} \left( T_u^V(x,y) - T_u(x-y) \right) u^{h-1} \, dy \, du \right. \\
        &  + \left. \int_{\rho^2(x)}^\infty \int_{\mathbb{R}^d \setminus B(x,1)} T_u^V(x,y) u^{h-1} \, dy \, du
        + \int_{\rho^2(x)}^\infty \int_{B(x,1)} T_u^V(x,y) u^{h-1} \, dy \, du
        + \frac{\rho(x)^{2h}}{h} \right) \mathcal{L}_V^{-(t-h)} f(x),
    \end{split}
\end{equation*}
where $0 < h < t$.


Since $\mathcal{L}_V^{-(t-h)} f\in \textup{Lip}_V^{\theta+t+h}$ we deduce that there exists $C>0$ such that
\begin{equation*}
    |\mathcal{L}_V^{-(t-h)} f(x) - \mathcal{L}_V^{-(t-h)} f(y)| \leq C |x-y|^{\theta + t - h}, \quad y\in \RR^d \text{ and } 0<h<t/2,
\end{equation*}
and
\begin{equation*}
    |\mathcal{L}_V^{-(t-h)} f(y)| \leq C \rho(y)^{\theta+t-h}, \quad y\in \RR^d \text{ and } 0<h<t/2.
\end{equation*}

Since $0<t<1-\theta$ we have that, for every $y\in \RR^d$ and $0<h<t/2$
\begin{equation*}
    \rho(y)^{\theta+t-h} \leq \rho(y)^{\theta +t/2} + \rho(y)^{\theta+t} \leq \rho(y)^\theta + \rho(y),
\end{equation*}
and 
\begin{equation*}
    |x-y|^{\theta +t -h } \leq |x-y|^{\theta +t} + |x-y|^{\theta+t/2} \leq |x-y|^\theta + |x-y|.
\end{equation*}

By taking into acccount that
\begin{equation*}
    \int_{\RR^d} \frac{\rho(y)^\theta + \rho(y)}{(1+|y|)^{d+2h_0}} dy < \infty 
\end{equation*}
for some $h_0\in(0,1)$,
we can use the Dominated Convergence Theorem to get
\begin{equation*}
    \begin{split}
        \lim_{h\rightarrow 0^+} \frac{\mathcal{L}_V^{-h} - I}{-h}\mathcal{L}_V^{-(t-h)} f(x) & = \lim_{h\rightarrow 0^+} \frac{1}{h} \left( \frac{1}{\Gamma(h)} \left( J_1(h) + J_2(h) + J_3(h) \right) - \mathcal{L}_V^{-t} f(x)\right)
        \\ & = -(\log \mathcal{L}_V) \mathcal{L}_V^{-t} f(x).
    \end{split}
\end{equation*}
Thus,~\eqref{eq: 3.3} is proved.

By combining~\eqref{eq: 3.1} and~\eqref{eq: 3.3} we conclude that $u(x,t) = \mathcal{L}_V^{-t} f(x)$, $x\in \RR^d$ and $t>0$, is in $D(\log \mathcal{L}_V)$ and 
\begin{equation*}
    \frac{\partial}{\partial t} u(x,t) = -\log(\mathcal{L}_V)) u(x,y),\quad x\in \RR^d \text{ and } t>0.
\end{equation*}

Since $\mathcal{L}_V^{-t}(f)\in \textup{Lip}_V^{\theta +t}$, $0<t<1-\theta$, the function $u(\cdot, t)$ is continuous in $\RR^d$ for every $t\in (0,1-\theta)$. 

In the first part of the proof we have proved that the function $u(x,\cdot)$ is continuous in $t\in(0,1-\theta)$, for every $x\in \RR^d$. We are going to see that the function $u$ is continuous in $\RR^d\times (0,1-\theta)$. Let $(x_0,t_0)\in \RR^d\times (0,1-\theta)$.  Suppose that, for every $k\in \NN$, $(x_k,t_k)\in \RR^d\times (0,1-\theta)$ and that $(x_k,t_k)\rightarrow (x_0,t_0)$, as $k\rightarrow \infty$. We can write
\begin{equation*}
    |u(x_k,t_k) - u(x_0,t_0)| \leq 
    | \mathcal{L}_V^{-t_k}(f)(x_k) - \mathcal{L}_V^{-t_k}(f)(x_0)| + |\mathcal{L}_V^{-t_k}(f)(x_0) - \mathcal{L}_V^{-t_0}(f)(x_0)|, \quad k\in\NN.
\end{equation*}

We chose $\delta_0>0$ such that $[t_0 - \delta_0,t_0+ \delta_0]\subseteq (0,1-\theta)$. There exists $k_0\in \NN$ such that $t_k \in [t_0-\delta_0, t_0+\delta_0]$ and $|x_k-x_0|<1$, for $k\geq k_0$. According to~\eqref{eq: 3.7}, there exists $C>0$ such that
\begin{equation}\label{eq: 3.8}
    |\mathcal{L}_V^{-t_k}(f)(x_k) - \mathcal{L}_V^{-t_{k}} (f)(x_0) | \leq C |x_k - x_0|^{\theta + t_k} \leq C |x_k - x_0|^{\theta+t_0-\delta_0}, \quad k\geq k_0.
\end{equation}

On the other hand, using~\eqref{eq: 1.2} we obtain 
\begin{equation*}
    \begin{split}
        |\mathcal{L}_V^{-t_k}(f)(x_0) - \mathcal{L}_V^{-t_0}(f)(x_0)| & \leq \int_{0}^\infty \left| \frac{u^{t_k}-1}{\Gamma(t_k)} - \frac{u^{t_0-1}}{\Gamma(t_0)}\right| \int_{\RR^d} T_u^V(x_0,y)|f(y)| dy du
        \\ & \leq \int_{0}^\infty \left| \frac{u^{t_k}-1}{\Gamma(t_k)} - \frac{u^{t_0-1}}{\Gamma(t_0)}\right| \int_{\RR^d} \frac{e^{\frac{-c|x_0-y|^2}{u}}}{u^{d/2}}|f(y)| dy du, \quad k\in \NN.
    \end{split}
\end{equation*}

Since $t_0+\delta_0-d/2<0$, we have that 
\begin{equation}\label{eq: 3.9}
\begin{split}
    \int_1^\infty \left| \frac{u^{t_k - 1}}{\Gamma(t_k)} - \frac{u^{t_0 - 1}}{\Gamma(t_0)} \right| \int_{\mathbb{R}^d} \frac{e^{-c|x_0 -y|^2/u}}{u^{d/2}} |f(y)| dy du
    &  \leq C \int_1^{\infty} \left( \frac{u^{t_k - 1}}{\Gamma(t_k)} + \frac{u^{t_0 - 1}}{\Gamma(t_0)} \right) \frac{1}{u^{d/2}} \int_{\RR^d} |f(y)| dy du
    \\ & \leq C \int_1^\infty u^{t_0+\delta_0-1-d/2} \int_{\RR^d} |f(y)| dy du <\infty, 
\end{split}
\end{equation}
for $k\geq k_0$. Since $\rho(y)\simeq \rho(x)$ when $|x_0 - y| < \rho(x)$ we get

\begin{equation}\label{eq: 3.10}
    \begin{split}
\int_0^1 & \left| \frac{u^{t_k-1}}{\Gamma(t_k)} - \frac{u^{t_0-1}}{\Gamma(t_0)}\right|  \int_{\RR^d} \frac{e^{-c|x_0-y|^2/u}}{u^{d/2}} |f(y)| dy du
 \\ & C\leq \int_0^1 u^{t_0 - \delta_0 -1} \left( \int_{|x_0 - y|<\rho(x_0)}  \frac{e^{-c|x_0-y|^2/u}}{u^{d/2}} |f(y)| dy   + \int_{|x_0 - y|\geq\rho(x_0)}  \frac{e^{-c|x_0-y|^2/u}}{u^{d/2}} |f(y)| dy 
\right) du
\\ & C\leq \int_0^1 u^{t_0 - \delta_0 -1} \left( \int_{|x_0 - y|<\rho(x_0)}  \frac{e^{-c|x_0-y|^2/u}}{u^{d/2}} \rho(y)^\theta dy   + \int_{|x_0 - y|\geq\rho(x_0)}  \frac{|f(y)|}{|x_0 - y|^d} dy 
\right) du
\\ & C\leq \int_0^1 u^{t_0 - \delta_0 -1} \left( \int_{\RR^d}  |f(y)| dy \rho(x_0)^{-d} + \rho(x_0)^\theta  
\right) du <\infty,
\end{split}
\end{equation}
for $k<k_0$. Nothe that the constants $C$ in~\eqref{eq: 3.9} and~\eqref{eq: 3.10} does not depend on $k$. 
By combining~\eqref{eq: 3.9} and~\eqref{eq: 3.10} and by using the Dominated Convergence Theorem we obtain
\begin{equation}\label{eq: 3.11}
    \lim_{k\rightarrow \infty} \mathcal{L}_V^{-t_k} (f)(x_0) = \mathcal{L}_V^{-t_0} (f)(x_0).
\end{equation}

By putting together~\eqref{eq: 3.8} and~\eqref{eq: 3.11} we conclude that

\begin{equation}
    \lim_{k\rightarrow \infty} \mathcal{L}_V^{-t_k} (f)(x_k) = \mathcal{L}_V^{-t_0} (f)(x_0).
\end{equation}
Thus, we prove that $u$ is continuous in $\RR^d \times (0,1-\theta)$. 

We now see that $\frac{\partial u}{\partial t}$ is continuous in $\RR^d \times (0,1-\theta)$. We have that

\begin{equation*}
    \frac{\partial u (x,t)}{\partial t} = \frac{1}{\Gamma(t)} \int_0^\infty u^{t-1} \log u \, T_u^V(f)(x) du {-} \frac{\Gamma'(t)}{\Gamma(t)} u(x,t), \qquad x\in\RR^d \text{ and } t\in (0,1-\delta). 
\end{equation*}

Differentiation under the integral sign is justified because we have that
\begin{equation*}
    \begin{split}
        \int_0^\infty u^{t-1} |\log u| \int_{\RR^d} T_u^V(x,y) |f(y)| dy du & \leq C \left( \int_0^1 u^{t-1} |\log u|  \int_{|x-y|<\rho(x)} \frac{e^{-c|x-y|^2/u}}{u^{d/2}} |f(y)| dy du \right.
        \\ & \qquad + \int_0^1 u^{t-1} |\log u|  \int_{|x-y|\geq\rho(x)} \frac{e^{-c|x-y|^2/u}}{u^{d/2}} |f(y)| dy du
        \\ & \qquad \left. + 
        \int_1^\infty u^{t-1} |\log u|  \int_{\RR^d} \frac{e^{-c|x-y|^2/u}}{u^{d/2}} |f(y)| dy du\right)
        \\ & \leq C  \left( \rho(x)^{-\theta}\int_0^1 u^{t-1} |\log u|  du  \right.
        \\ & \qquad +  \rho(x)^{-d} \int_0^1 u^{t-1} |\log u|   \int_{|x-y|\geq \rho(x)} {|f(y)| dy du }
        \\ & \qquad \left. + 
        \int_1^\infty u^{t-1-d/2} |\log u|  \int_{\RR^d}  |f(y)| dy du\right) <\infty,
    \end{split}
\end{equation*}
for $x\in \RR^d$ and $t_0\in(0,1-\theta)$. We choose a sequence $\{(x_k,t_k)\}_{k\in \NN}\subseteq \RR^d \times (0,1-\theta)$ such that  $\displaystyle \lim_{k\rightarrow \infty} (x_k,t_k) = (x_0,t_0)$. We can write
\begin{equation*}
\begin{split}
    \int_0^\infty u^{t_k-1} \log u T_u^V(f)(x_k)du & = \int_1^\infty u^{t_k-1} \log u \int_{\RR^d} T_u^V (x_k,y) f(y) dy du 
    \\ & \qquad +
    \int_{0}^1 u^{t_k -1} \log u \int_{|y-x_0|\geq \rho(x_0)} T_u^V (x_k,y) f(y) dy du
    \\ & \qquad + \int_{0}^1 u^{t_k -1} \log u \int_{|y-x_0|< \rho(x_0)} T_u^V (x_k,y) f(y) dy du
    \\ & = J_1(k) + J_2(k) + J_3(k), \qquad  k\in \NN.
\end{split}
\end{equation*}

We choose $\delta_0>0$ such that $[t_0 - \delta_0, t_0 + \delta_0]\subseteq (0,1-\theta)$. There exists $k_0\in \NN$ such that $t_k \in (t_0 - \delta_0, t_0 + \delta_0)$, $k\geq k_0$. 

We now apply~\eqref{eq: 1.2} to get
\begin{equation*}
    \begin{split}
        u^{t_k-1} |\log u | T_u^V(x_k,y) |f(y)| \leq C u^{t_0+\delta_0 - 1} |\log u | \frac{e^{-c|y-x_k|^2/u}}{u^{d/2}} |f(y)| \leq C |\log u | u^{{t_0 + \delta_0 -d/2- 1}} |f(y)|, 
    \end{split}
\end{equation*}
for $y\in \RR^d$, $u\in(1,\infty)$ and $k\geq k_0$, being
\[
\int_1^\infty \int_{\RR^d} u^{t_0+\delta_0 - 1 - d/2} |f(y)|dy du <\infty. 
\]
Then
\[
\lim_{k\rightarrow\infty} J_1(k) = \int_1^\infty u^{t_0-1} \log u \int_{\RR^d} T_u^V(x_0,y) f(y) dy.
\]

On the other hand, there exists $k_1\geq k_0$, $k_1\in \NN$ such that $|x_0 - x_k| < \rho(x_0)/2$. We obtain
\begin{equation*}
    \begin{split}
        u^{t_k-1}|\log u | T_u^V(x_k,y) |f(y)|  & \leq C u^{t_0 - \delta_0 - 1} |\log u| \frac{e^{-c|y-x_k|^2/u}}{u^{d/2}} |f(y)|
        \\ & \leq  C u^{t_0 - \delta_0 - 1} |\log u| \frac{|f(y)|}{|y-x^k|^d}
        \\ & \leq C u^{t_0 - \delta_0 - 1} |\log u| |f(y)| \rho(x_0)^{-d}, 
    \end{split}
\end{equation*}
for $y\in \RR^d$, $|y-x_0|\geq \rho(x_0)$,  $0<t<1$ and $k\geq k_1$, because $|y-x_k| \geq |y-x_0| - |x_0 - x_k| \geq \rho(x_0) - \rho(x_0)/2 = \rho(x_0)/2$, when $|y-x_0|\geq \rho(x_0)$ and $k\geq k_1$. 

Since 
\[\int_0^1 \int_{\RR^d} u^{t_0-\delta_0 - 1} |\log u| |f(y)|dy du <\infty,\]
it follows that
\[\lim_{k\rightarrow \infty} J_2(k) = \int_0^1 u^{t_0-1} \log u \int_{{|x-y|\geq \rho(x_0)}} T_u^V(x_0,y) f(y) dy.\]


We also have that, since $f\in \Lip_V^\theta$, we get
\begin{equation*}
    \begin{split}
        u^{t_k-1} |\log u | \int_{|y-x_0|<\rho(x_0)} T_t(x_k,y) |f(y)| dy & \leq C u^{t_0 - {\delta_0} - 1} |\log u| \int_{|y-x_0|<\rho(x_0)} \frac{e^{-c|x_k-y|^2/u|}}{u^{d/2}} |f(y)| dy 
        \\ & \leq C u^{t_0 - {\delta_0} - 1} |\log u| \int_{|y-x_0|<\rho(x_0)} \frac{e^{-c|x_k-y|^2/u|}}{u^{d/2}} \rho(y)^\theta dy
        \\ & \leq C \rho(x_0)^\theta  u^{t_0 - {\delta_0} - 1} |\log u| \int_{|y-x_0|<\rho(x_0)} \frac{e^{-c|x_k-y|^2/u|}}{u^{d/2}} dy,
    \end{split}
\end{equation*}
for $u\in (0,1)$ and $k\geq k_0$, and 
\[ T_u^V(x_k,y) |f(y)|  \leq C \frac{e^{-c|x_k-y|^2/2}}{u^{d/2}}|f(y)| \leq C\frac{|f(y)|}{u^{d/2}},\]
for $u>0$, $y\in \RR^d$ and $k\in \NN$.
Then
\[\lim_{k\rightarrow \infty } \int_{|y-x_0|<\rho(x_0)} T_u^V(x_k,y) f(y) dy = \int_{|y-x_0|<\rho(x_0)} T_u^V(x_0,y) f(y) dy, \qquad u>0,\]
and
\[
\begin{split}
\lim_{k\rightarrow \infty } J_3(k) & = \int_0^1 u^{t_0-1} \log u \lim_{k\rightarrow \infty }  \int_{|y-x_0|<\rho(x_0)} T_u^V(x_k,y) f(y) dy du
\\ & = \int_0^1 u^{t_0-1} \log u  \int_{|y-x_0|<\rho(x_0)} T_u^V(x_0,y) f(y) dy du
\end{split}
\]

Then, we prove that 
\[\lim_{k\rightarrow\infty} \int_0^\infty u^{t_k-1} \log u   T_u^V(f)(x_k) du = \int_0^\infty u^{t_0-1} \log u   T_u^V(f)(x_0) du.\]
We conclude that $\frac{\partial u}{\partial t}$, and hence $(\log \mathcal{L}_V)u$ is continuous in $\RR^d\times (0,1-\theta)$.


Our next objective is to see that
\begin{equation}\label{eq: 3.12}
\lim_{t\rightarrow 0^+} u(x,t) = f(x), \qquad x\in \RR^d.    
\end{equation}
Let $x\in \RR^d$. We choose $R>0$ such that $R>|x|+1$. We decompose $u$ as follows.
\begin{equation*}
    \begin{split}
        u(x,t) & = \frac{1}{\Gamma(t)} \int_0^\infty u^{t-1} \int_{B(x,R)} T_u^V(x,y) f(y) dy du + \frac{1}{\Gamma(t)} \int_0^\infty u^{t-1} \int_{\RR^d\setminus B(x,R)} T_u^V(x,y) f(y) dy du
        \\ & = J_1(t) + J_2(t), \qquad 0<t<1-\theta.
    \end{split}
\end{equation*}

First we prove that $J_1(t) \rightarrow f(x)$ as $t\rightarrow 0^+$. We can write
\begin{equation*}
    \begin{split}
        J_1(t) & = \frac{1}{\Gamma(t)} \int_0^{\rho^2(x)} u^{t-1} \int_{B(x,R)} \left( T_u^V(x,y) - T_u(x-y)\right) f(y) dy
        \\ & \qquad + \frac{1}{\Gamma(t)} \int_{\rho^2(x)}^\infty u^{t-1} \int_{B(x,R)}  T_u^V(x,y)  f(y) dy
        \\ & \qquad - \frac{1}{\Gamma(t)} \int_{\rho^2(x)}^\infty u^{t-1} \int_{B(x,R)}  T_u(x-y)  f(y) dy
        \\ & \qquad + \frac{1}{\Gamma(t)} \int_{0}^\infty u^{t-1} \int_{B(x,R)}  T_u(x-y)  f(y) dy
        \\ & = J_{1,1}(t) +J_{1,2}(t) + J_{1,3}(t) + J_{1,4}(t), \qquad 0<t<1-\theta.
    \end{split}
\end{equation*}
According to~\eqref{eq: 1.5} we have that
\[ |J_{1,1}(t)| \leq C \frac{1}{\Gamma(t)}
\int_0^{\rho^2(x)} u^{t-1} \int_{B(x,R)} \left(\frac{\sqrt{u}}{\rho(x)} \right)^\delta \omega_k(x-y) |f(y)| dy du, \qquad 0<t<1-\theta.\]

By using~\eqref{eq: rox_vs_roy} and a compactness argument as we did in the proof of Theorem~\ref{thm: 1.2}\ref{item: teo 1.2 - b}, we can find $C>1$ such that
\[\frac{1}{C} \leq \rho(y)\leq C, \qquad y\in B(x,R).\]

Since $f\in \Lip_V^\theta$ we deduce that
\[
|J_{1,1}(t)| \leq \frac{C}{\Gamma(t)} \int_0^{\rho^2(x)} u^{t-1+\delta/2} \int_{B(x,R)} \omega_u(x-y) dy du \leq \frac{C}{\Gamma(t)} \frac{\rho(x)^{2(t+\delta/2)}}{t+\delta/2}, \qquad 0<t<1-\theta.
\]
Then, $\displaystyle\lim_{t\rightarrow 0^+} J_{1,1}(t) = 0$. 

From \eqref{eq: 1.2} we get
\begin{align*}
    \left|J_{1,2}(t)\right|+ \left|J_{1,3}(t)\right| &\leq  C\frac{1}{\Gamma(t)}\int_{\rho^2(x)}^\infty u^{t-1}\int_{B(x,R}\frac{e^{-c\frac{|x-y|^2}{u}}}{u^{\frac{d}{2}}}|f(y)|dydu\\
    & \leq  C\frac{1}{\Gamma(t)} \int_{\rho^2(x)}^\infty u^{t- \frac{\delta}{2}-1}du \int_{\mathbb R^d} |f(y)|dy \\
    & \leq  C \frac{1}{\Gamma(t)}\frac{\rho(x)^{2t-n}}{\frac{d}{2}-t},\quad 0<t<1-\theta 
\end{align*}
It follows that $\displaystyle\lim_{t\rightarrow 0}J_{1,2}(t)+J_{1,3}(t))=0$.

By interchanging the order of integration we obtain
\begin{align*}
    J_{1,4}(t) &=\frac{1}{\Gamma(t)}\int_{B(x,R)}f(y)\int_0^\infty\frac{e^{-\frac{|x-y|^2}{4u}}}{(4\pi u)^{\frac{d}{2}}}u^{t-1}du dy\\
    &= \frac{4^{-t}\Gamma\left(\frac{d-2t}{2}\right)}{\pi^{\frac{d}{2}\Gamma(t)}}\int_{B(x,R)}\frac{f(y)}{|x-y|^{d-2t}}dy,\quad 0<t<1-\theta.
\end{align*}
As it was proved in the proof of \cite[Proposition 2.3]{ChV1} we have that
\[
\lim_{t \rightarrow 0} J_{1,4}(t) = f(x).
\]
On the other hand by using again \eqref{eq: 1.2} we get
\begin{align*}
    J_{2}(t) &=\frac{1}{\Gamma(t)}\int_0^\infty u^{t-1}\int_{\mathbb R^d\setminus B(x,R)}\frac{e^{-\frac{|x-y|^2}{4u}}}{(4\pi u)^{\frac{n}{2}}}|f(y)|dy du\\ 
    &= \frac{4^{-t}\Gamma\left(\frac{d-2t}{2}\right)}{\pi^{\frac{d}{2}\Gamma(t)}}\int_{\mathbb R^d \setminus B(x,R)}\frac{|f(y)|}{|x-y|^{n-2t}}dy,\quad 0<t<1-\theta.
\end{align*}

By using now H\"older and Gagliardo-Nirenberg inequalities as in the proof of \cite[Proposition 2.3]{ChV1} we conclude that
\[
\lim_{t \rightarrow 0} J_2(t)=0.
\]



\bibliographystyle{acm}


\begin{thebibliography}{10}

\bibitem{BFHR}
{\sc Betancor, J.~J., Fari\~{n}a, J.~C., Harboure, E., and Rodr\'{\i}guez-Mesa, L.}
\newblock Variation operators for semigroups and {R}iesz transforms on {$BMO$} in the {S}chr\"{o}dinger setting.
\newblock {\em Potential Anal. 38}, 3 (2013), 711--739.

\bibitem{BHQ}
{\sc Bongioanni, B., Harboure, E., and Quijano, P.}
\newblock Weighted inequalities for {S}chr\"{o}dinger type singular integrals.
\newblock {\em J. Fourier Anal. Appl. 25}, 3 (2019), 595--632.

\bibitem{BHS}
{\sc Bongioanni, B., Harboure, E., and Salinas, O.}
\newblock Weighted inequalities for negative powers of {S}chr\"{o}dinger operators.
\newblock {\em J. Math. Anal. Appl. 348}, 1 (2008), 12--27.

\bibitem{CS}
{\sc Caffarelli, L., and Silvestre, L.}
\newblock An extension problem related to the fractional {L}aplacian.
\newblock {\em Comm. Partial Differential Equations 32}, 7-9 (2007), 1245--1260.

\bibitem{Ch-orderm}
{\sc Chen, H.}
\newblock On {$m$}-order logarithmic {L}aplacians and the applications.
\newblock {\em Anal. Appl. (Singap.) 24}, 2 (2026), 419--461.

\bibitem{ChHW}
{\sc Chen, H., Hauer, D., and Weth, T.}
\newblock An extension problem for the logarithmic {L}aplacian, 2023.
\newblock \href{https://arxiv.org/abs/2312.15689}{arXiv.2312.15689}.

\bibitem{ChV2}
{\sc Chen, H., and V\'{e}ron, L.}
\newblock Bounds for eigenvalues of the {D}irichlet problem for the logarithmic {L}aplacian.
\newblock {\em Adv. Calc. Var. 16}, 3 (2023), 541--558.

\bibitem{ChV1}
{\sc Chen, H., and V\'{e}ron, L.}
\newblock The {C}auchy problem associated to the logarithmic {L}aplacian with an application to the fundamental solution.
\newblock {\em J. Funct. Anal. 287}, 3 (2024), Paper No. 110470, 72.

\bibitem{ChW}
{\sc Chen, H., and Weth, T.}
\newblock The {D}irichlet problem for the logarithmic {L}aplacian.
\newblock {\em Comm. Partial Differential Equations 44}, 11 (2019), 1100--1139.

\bibitem{Ch-manifolds}
{\sc Chen, R.}
\newblock Logarithmic {L}aplacian on {G}eneral {R}iemannian {M}anifolds, 2025.
\newblock \href{https://arxiv.org/abs/2506.19311}{arXiv.2506.19311}.

\bibitem{ChX}
{\sc Chen, R., and Xu, W.}
\newblock The {L}ogarithmic {L}aplacian on {G}eneral {G}raphs, 2025.
\newblock \href{https://arxiv.org/abs/2507.05936}{arXiv.2507.05936}.

\bibitem{CF}
{\sc Coifman, R.~R., and Fefferman, C.}
\newblock Weighted norm inequalities for maximal functions and singular integrals.
\newblock {\em Studia Math. 51\/} (1974), 241--250.

\bibitem{dLT}
{\sc De~Le\'{o}n-Contreras, M., and Torrea, J.~L.}
\newblock Lipschitz spaces adapted to {S}chr\"{o}dinger operators and regularity properties.
\newblock {\em Rev. Mat. Complut. 34}, 2 (2021), 357--388.

\bibitem{Dz1}
{\sc Dziuba\'{n}ski, J.}
\newblock Spectral multiplier theorem for {$H^1$} spaces associated with some {S}chr\"{o}dinger operators.
\newblock {\em Proc. Amer. Math. Soc. 127}, 12 (1999), 3605--3613.

\bibitem{Dz2}
{\sc Dziuba\'{n}ski, J.}
\newblock A spectral multiplier theorem for {$H^1$} spaces associated with {S}chr\"{o}dinger operators with potentials satisfying a reverse {H}\"{o}lder inequality.
\newblock {\em Illinois J. Math. 45}, 4 (2001), 1301--1313.

\bibitem{DGMTZ}
{\sc Dziuba\'{n}ski, J., Garrig\'{o}s, G., Mart\'{\i}nez, T., Torrea, J.~L., and Zienkiewicz, J.}
\newblock {$BMO$} spaces related to {S}chr\"{o}dinger operators with potentials satisfying a reverse {H}\"{o}lder inequality.
\newblock {\em Math. Z. 249}, 2 (2005), 329--356.

\bibitem{DzZ3}
{\sc Dziuba\'{n}ski, J., and Zienkiewicz, J.}
\newblock {$H^p$} spaces for {S}chr\"{o}dinger operators.
\newblock In {\em Fourier analysis and related topics ({B}\k{e}dlewo, 2000)}, vol.~56 of {\em Banach Center Publ.} Polish Acad. Sci. Inst. Math., Warsaw, 2002, pp.~45--53.

\bibitem{DzZ2}
{\sc Dziuba\'{n}ski, J., and Zienkiewicz, J.}
\newblock {$H^p$} spaces associated with {S}chr\"{o}dinger operators with potentials from reverse {H}\"{o}lder classes.
\newblock {\em Colloq. Math. 98}, 1 (2003), 5--38.

\bibitem{FaFe2}
{\sc Fall, M.~M., and Felli, V.}
\newblock Sharp essential self-adjointness of relativistic {S}chr\"{o}dinger operators with a singular potential.
\newblock {\em J. Funct. Anal. 267}, 6 (2014), 1851--1877.

\bibitem{FaFe1}
{\sc Fall, M.~M., and Felli, V.}
\newblock Unique continuation properties for relativistic {S}chr\"{o}dinger operators with a singular potential.
\newblock {\em Discrete Contin. Dyn. Syst. 35}, 12 (2015), 5827--5867.

\bibitem{Feu1}
{\sc Feulefack, P.~A.}
\newblock The logarithmic {S}chr\"{o}dinger operator and associated {D}irichlet problems.
\newblock {\em J. Math. Anal. Appl. 517}, 2 (2023), Paper No. 126656, 33.

\bibitem{Feu2}
{\sc Feulefack, P.~A.}
\newblock The fractional logarithmic {S}chr\"{o}dinger operator: properties and functional spaces.
\newblock {\em J. Pseudo-Differ. Oper. Appl. 15}, 3 (2024), Paper No. 55, 38.

\bibitem{FJW}
{\sc Feulefack, P.~A., Jarohs, S., and Weth, T.}
\newblock Small order asymptotics of the {D}irichlet eigenvalue problem for the fractional {L}aplacian.
\newblock {\em J. Fourier Anal. Appl. 28}, 2 (2022), Paper No. 18, 44.

\bibitem{GHY}
{\sc Ge, A., He, Q., and Yan, D.}
\newblock On weighted compactness of oscillation and variation of commutators associated with {S}chr\"{o}dinger operators.
\newblock {\em Anal. Math. 49}, 3 (2023), 765--805.

\bibitem{G}
{\sc Gehring, F.~W.}
\newblock The {$L\sp{p}$}-integrability of the partial derivatives of a quasiconformal mapping.
\newblock {\em Acta Math. 130\/} (1973), 265--277.

\bibitem{GR}
{\sc Gradshteyn, I.~S., and Ryzhik, I.~M.}
\newblock {\em Table of integrals, series, and products}, eighth~ed.
\newblock Elsevier/Academic Press, Amsterdam, 2015.
\newblock Translated from the Russian, Translation edited and with a preface by Daniel Zwillinger and Victor Moll, Revised from the seventh edition [MR2360010].

\bibitem{HLW}
{\sc Harrach, B., Lin, Y.-H., and Weth, T.}
\newblock The {C}alder\'{o}n problem for the logarithmic {S}chr\"{o}dinger equation.
\newblock {\em J. Differential Equations 444\/} (2025), Paper No. 113665, 25.

\bibitem{HLS}
{\sc Hern\'{a}ndez-Santamar\'{\i}a, V., L\'{o}pez~R\'{\i}os, L.~F., and Salda\~{n}a, A.}
\newblock Optimal boundary regularity and a {H}opf-type lemma for {D}irichlet problems involving the logarithmic {L}aplacian.
\newblock {\em Discrete Contin. Dyn. Syst. 45}, 1 (2025), 1--36.

\bibitem{JX}
{\sc Jin, T., and Xiong, J.}
\newblock A fractional {Y}amabe flow and some applications.
\newblock {\em J. Reine Angew. Math. 696\/} (2014), 187--223.

\bibitem{Kur}
{\sc Kurata, K.}
\newblock An estimate on the heat kernel of magnetic {S}chr\"{o}dinger operators and uniformly elliptic operators with non-negative potentials.
\newblock {\em J. London Math. Soc. (2) 62}, 3 (2000), 885--903.

\bibitem{LW}
{\sc Laptev, A., and Weth, T.}
\newblock Spectral properties of the logarithmic {L}aplacian.
\newblock {\em Anal. Math. Phys. 11}, 3 (2021), Paper No. 133, 24.

\bibitem{Lee}
{\sc Lee, D.}
\newblock Fundamental {S}olutions of the {L}ogarithmic {L}aplacian: {A}n {A}pproach via the {D}ivision {P}roblem, 2025.
\newblock \href{https://arxiv.org/abs/2506.20121}{arXiv.2506.20121}.

\bibitem{LL}
{\sc Li, C., and L\"{u}, Y.}
\newblock Maximum principles for {L}aplacian and fractional {L}aplacian with critical integrability.
\newblock {\em J. Geom. Anal. 33}, 7 (2023), Paper No. 203, 27.

\bibitem{LPW}
{\sc Li, Q., Peng, S., and Wen, S.}
\newblock Existence of the least energy sign-changing solutions for fractional {B}rezis-{N}irenberg problem.
\newblock {\em Adv. Differential Equations 30}, 1-2 (2025), 69--92.

\bibitem{Lu}
{\sc Ludwig, M.}
\newblock Anisotropic fractional perimeters.
\newblock {\em J. Differential Geom. 96}, 1 (2014), 77--93.

\bibitem{MSTZ}
{\sc Ma, T., Stinga, P.~R., Torrea, J.~L., and Zhang, C.}
\newblock Regularity properties of {S}chr\"{o}dinger operators.
\newblock {\em J. Math. Anal. Appl. 388}, 2 (2012), 817--837.

\bibitem{ReSi}
{\sc Reed, M., and Simon, B.}
\newblock {\em Methods of modern mathematical physics. {I}. {F}unctional analysis}.
\newblock Academic Press, New York-London, 1972.

\bibitem{RS}
{\sc Ros-Oton, X., and Serra, J.}
\newblock The {D}irichlet problem for the fractional {L}aplacian: regularity up to the boundary.
\newblock {\em J. Math. Pures Appl. (9) 101}, 3 (2014), 275--302.

\bibitem{SV}
{\sc Servadei, R., and Valdinoci, E.}
\newblock The {B}rezis-{N}irenberg result for the fractional {L}aplacian.
\newblock {\em Trans. Amer. Math. Soc. 367}, 1 (2015), 67--102.

\bibitem{Shen94}
{\sc Shen, Z.~W.}
\newblock On the {N}eumann problem for {S}chr\"{o}dinger operators in {L}ipschitz domains.
\newblock {\em Indiana Univ. Math. J. 43}, 1 (1994), 143--176.

\bibitem{Shen95}
{\sc Shen, Z.~W.}
\newblock {$L^p$} estimates for {S}chr\"{o}dinger operators with certain potentials.
\newblock {\em Ann. Inst. Fourier (Grenoble) 45}, 2 (1995), 513--546.

\bibitem{Yo}
{\sc Yosida, K.}
\newblock {\em Functional analysis}, sixth~ed., vol.~123 of {\em Grundlehren der Mathematischen Wissenschaften}.
\newblock Springer-Verlag, Berlin-New York, 1980.

\bibitem{ZT}
{\sc Zhang, Q., and Tang, L.}
\newblock Variation operators on weighted {H}ardy and {BMO} spaces in the {S}chr\"{o}dinger setting.
\newblock {\em Bull. Malays. Math. Sci. Soc. 45}, 5 (2022), 2285--2312.

\end{thebibliography}



\end{document}
